% =====================================================================
%  Harald Høffding, FORMEL LOGIK.  TIL BRUG VED FORELÆSNINGER.  Femte omarbejdede Udgave.
%  Kjøbenhavn og Kristiania: Gyldendalske Boghandel, Nordisk Forlag, 1907 (Forlagets Bogtrykkeri).  [X] + 48 pp.  Antiqua.
%  The edition Høffding sent in sheets to the student Niels Bohr for criticism in November 1906 (Pais, Niels Bohr's Times,
%  ch. 6 nn. 45-47); the Forord (December 1906) thanks "en af mine tidligere Tilhørere" and Dr. Anton Thomsen.
%
%  GENERATED; DO NOT EDIT BY HAND.  Built by ebook-method/formel-logik/rebuild.sh (draft.py --emit, collate/check.py --apply,
%  draft.py --post, heads.py, notatex.py); edit the pipeline (template.tex, patches.py, decisions.tsv, schemes.tsv,
%  casenums.tsv, italics.tsv), never this file.  Rebuild line in texts/hoeffding/formel-logik/RESUME-NOTES.md.
%
%  EDITION AND SCAN.  Text of record: the University of Minnesota copy, scanned by Google and served by HathiTrust
%  (umn.31951002000878s): ~/bibliotek/Høffding, Harald/1907-formel-logik-5udg.pdf (SCANS.tsv by sha256 b176ccbb…),
%  65 PDF pages, 600 ppi bitonal, Google's OCR layer.  PDF = printed + 13 (p. 1 = PDF 14, p. 48 = PDF 61); every folio
%  present.  PDF 6 half-title, 8 title page, 10 Forord til femte Udgave, 12 Indholdsfortegnelse (one page; PDF 13 blank).
%  The front pages carry no folio: their markers [III] [V] [VII] [IX] are counted back from p. 1.  No e-book and no earlier
%  transcription of this edition exists (edition check 2026-10-29); the first edition (1884) is a different text.
%
%  NOT TRANSCRIBED (marks of this copy): the bookplate (PDF 3) and a library stamp (PDF 11); on the title page the pencil
%  date "1843-" after the author's name and the library's perforation stamp; on the Forord the pencil shelfmark "160 H67",
%  the accession stamp "50811" and the dealer's stamp "Dec. 14. 1907. Stechert" (which overprints the line "Jeg har tillige
%  ... Bemærkninger fra en af mine tidligere Tilhørere": its words were read at the image); Google's watermark; the sheet
%  signatures ("Harald Høffding: Formel Logik. 5. Udg. N", "N*").
%
%  METHOD AND ACCURACY (2026-10-29; TRANSCRIPTION-PLAYBOOK.md §00).  No one re-typed the text.
%  (1) The draft is Google's layer turned into LaTeX by script (draft.py); corrections that a word decision cannot express
%      are literal patches (patches.py: 73 before the collation, 5 after it, 6 paragraph joins, 8 dashes moved).
%  (2) It was collated word by word, with punctuation, against an independent tesseract 4 reading (dan, tessdata_best,
%      300 dpi) by ebook-method/collate/check.py: 454 disagreements in pp. 1-48 at the end; 10 settled by check.py's rules,
%      444 by decisions -- 88 by rules.py (dashes the layer drops, quote kinds, accents; never on a line of notation),
%      252 from the blind notation readings (below), 10 by the three-way vote, 97 at the image in the main session -- OPEN 0.
%  (3) The LOGICAL NOTATION was read BLIND: every line with notation (296) and the 17 schemes and tables, as crops, by six
%      readers who saw neither machine reading (NOTATION-BRIEF.md; 12 planted control lines, 12 right).  Each reading was
%      aligned with the draft token by token (notareads.py) and every difference put to a three-way vote of reader, layer and
%      tesseract (vote.py): 60 places where reader and tesseract agreed against the draft were corrected; the residue (70)
%      was settled at the image in the main session (z-sheets), including the whole p. 36 table, the p. 40 line the layer
%      had scrambled, and the p. 45 scheme.  The schemes are set whole from the adjudicated readings (schemes.tsv).
%  (4) The 97 open prose checks were settled at the image in the main session (p-sheets); 8 words divided over a page turn,
%      whose hyphen the layer had lost on the page's last line, were joined.
%  (5) Read at the image in the main session: every heading (wording, case, bold, letterspacing), the front matter, the 14
%      paragraph-initial long dashes (an ink test, then the image: 8 end the paragraph before, 6 stand flush inside one),
%      the exclusion sign of p. 27.
%  (6) ITALICS: found by a stroke-slant test on tesseract's word boxes (calibrated on words the blind readers saw italic:
%      ITAL 1.06), merged into runs across short words (italics.py), and then EVERY run read at the image in the main session
%      (it-, it2-, it3-sheets): the test misses the first word(s) of italic sentences and stops at short words (extents set
%      in italics-fix.tsv), and gave 4 false positives (p. 2 "fast", p. 24 "Grund", p. 31 "(Naar …)", and one crop that
%      showed the wrong word), dropped; a final stop after an italic run is roman.  72 runs set.
%  (7) Sampled blind page check (Forord, pp. 13, 38, 44; 1,122 words; 2/2 planted controls found): 0 word errors,
%      1 punctuation (Uranus’), 5 italic sentence openings -- all mended; every italic run then verified at the image
%      (main session).
%  Cost: approximately 1.3M agent tokens (the pipeline agent ~0.73M by its own token count; six blind notation readers
%  ~490k; the blind page reader ~97k), plus the main session's adjudication at the crop sheets.
%
%  CONVENTIONS.  Diplomatic: spelling, capitals and punctuation as printed; „…“ quotation marks; the long dash as ---;
%  »ɔ:« (det er) as printed.  The paragraph numbers (1.-47.) as printed.  Words divided at a line end are joined; a word
%  divided across a page turn carries the marker inside it (\opage{N}).  Running heads are kept as comments at each page
%  turn.  LOGICAL NOTATION in math mode (the print sets the logical letters in italics): =, <, > as printed; the print's
%  negation sign ~ ("er ikke") as \sim; the exclusion sign ·|· of p. 27 as \udel; subscripts as _1; the raised case numbers
%  ¹) ²) ... as \textsuperscript{N}); schemes as display arrays, the printed rule as \hline, the short dotted rule of the
%  three-column schemes pp. 44-45 as \cdashline; the p. 36 table as two tabulars.  Headings: chapter numeral and title in
%  capitals; the section heads A.-D., a./b. and the Tillæg's 1./2. in bold (in a./b. the parenthesis roman); the α./β.
%  heads letterspaced (\spaced).  Italics \textit; names in small capitals (the Forord) \textsc.
%
%  THE CONTENTS PAGE AND THE HEADS DISAGREE (both recorded; each is set as printed):
%    p. 11  Indhold »C. Dommes Dannelse og Formulering«;  head »C. Dommens Dannelse og Formulering.«
%    p. 35  Indhold »α. Indirekte Slutning ved ren Deduktion«;  head »α. Indirekte Slutning som ren Deduktion.«
%    p. 42  Indhold »1. Mill’s Induktionsmetoder«;  head »1. Mill’s induktive Metoder.«
%
%  PRINTER'S ERRORS, PRINTED AS IS: p. 45, the scheme of the proportional variations: the middle cell of the bottom row reads
%  A₂V₁X₁ (B₂V₁X₁ expected?), and B₁V₁X has no subscript on X.  Logged at the site with % PRINTED AS IS.
%  After the translation (2026-10-29), the translators' 19 odd spots were read at the image: also PRINTED AS IS (logged at the
%  site): p. VII »pædagisk«, p. 3 and p. 34 a semicolon before a relative clause, p. 14 »af / af sig selv«, p. 17 »deres
%  Indhold forekomme«, p. 24 »Videnkaber«, p. 33 the cross-reference »(33)« (32 meant?).  Mended: 8 notation spans the notation
%  pass had missed or set wrongly (mathfix.tsv, casenums.tsv), words it had set in math mode (notatex.py), p. 36 »Varme- og«.
% =====================================================================
\documentclass[12pt,a4paper,openany]{book}
\usepackage[utf8]{inputenc}
\usepackage[T1]{fontenc}
\usepackage{libertinus}
\usepackage{libertinust1math}
\usepackage{array}
\usepackage{arydshln}   % formel-logik: \cdashline, the short dotted rule of the schemes pp. 44-45 (notatex.py)
\usepackage{graphicx}
\usepackage{newunicodechar}
\newunicodechar{ɔ}{\reflectbox{c}}   % formel-logik: »ɔ:« (det er), as in the Brøchner transcriptions
\newcommand{\udel}{\mathbin{\cdot\kern-0.1em|\kern-0.1em\cdot}}   % formel-logik: the exclusion sign ·|· of p. 27 (»gensidigt Udelukkelsesforhold«)
\usepackage{textalpha}
\usepackage[danish]{babel}
\usepackage[protrusion=true,expansion=true]{microtype}
\usepackage{setspace}
\usepackage[margin=1.3in]{geometry}
\usepackage{fancyhdr}
\usepackage{hyperref}
\hypersetup{pdftitle={Formel Logik}, pdfauthor={Harald Høffding}, hidelinks}
\pagestyle{fancy}
\fancyhf{}
\fancyhead[LE,RO]{\thepage}
\fancyhead[C]{\textit{Høffding: Formel Logik, 5. Udg. (1907)}}
\setlength{\headheight}{15pt}
\setstretch{1.3}
\usepackage{marginnote}
\newcommand{\opage}[1]{\marginnote{\footnotesize\textit{[#1]}}}
\newcommand{\whiteline}{\par\vspace{\baselineskip}\par}
% heading macros (from ledende-tanker; to be used by patch once the headings are read at the image)
\newcommand{\parthead}[1]{\par\bigskip\phantomsection\begin{center}{\Large #1}\end{center}\addcontentsline{toc}{part}{#1}\medskip}
\newcommand{\unnumhead}[1]{\par\bigskip\phantomsection\begin{center}{\Large #1}\end{center}\addcontentsline{toc}{chapter}{#1}\medskip}
\newcommand{\chapterhead}[2]{\par\bigskip\phantomsection\begin{center}{\large #1\ #2}\end{center}\addcontentsline{toc}{chapter}{#1 #2}\medskip}
\newcommand{\chapterheadii}[3]{\par\bigskip\phantomsection\begin{center}{\large #1\ #2\\ #3}\end{center}\addcontentsline{toc}{chapter}{#1 #2 #3}\medskip}
\newcommand{\sectionhead}[2]{\par\bigskip\phantomsection\begin{center}#1 \textit{#2}\end{center}\addcontentsline{toc}{section}{#1 #2}\medskip\ignorespaces}
\newcommand{\ix}[3]{\par\noindent\hspace*{#1}#2\dotfill #3}
% formel-logik (heads.py): chapter = numeral, title in capitals, the smaller subtitle; section heads; letterspacing
\newcommand{\chaphead}[3]{\par\bigskip\phantomsection\begin{center}\ifx&#1&\else{\large #1}\par\smallskip\fi{\Large #2}\ifx&#3&\else\par\smallskip{\large #3}\fi\end{center}\addcontentsline{toc}{chapter}{\ifx&#1&\else #1. \fi #2}\medskip}
\newcommand{\sechead}[2]{\par\bigskip\phantomsection\begin{center}#2\end{center}\addcontentsline{toc}{section}{#1}\medskip}
\newcommand{\spaced}[1]{\textls[180]{#1}}

\begin{document}

\frontmatter
% --- front matter (PDF 6 half-title, 8 title page, 10 Forord, 12 Indhold) as drafted from the layer: the title page
%     and the Indhold are to be set in their printed form by patch after they are read at the image ---
% --- p. III ---
\thispagestyle{empty}\opage{III}\vspace*{5cm}
\begin{center}{\Large FORMEL LOGIK}\end{center}
\clearpage

% --- p. V ---
% the title page (PDF 8); not transcribed: the pencil date "1843-" after the name and the library's perforation stamp
\begin{titlepage}\centering\opage{V}
\vspace*{1cm}
{\large HARALD HØFFDING}\par\medskip
\rule{3cm}{0.4pt}\par\vspace{1.5cm}
{\Huge FORMEL LOGIK}\par\bigskip
{\large TIL BRUG VED FORELÆSNINGER}\par\vspace{1.2cm}
\textit{FEMTE OMARBEJDEDE UDGAVE}\par
\vspace{2cm}
% [the publisher's device]
\vspace{3cm}
KJØBENHAVN OG KRISTIANIA\par
GYLDENDALSKE BOGHANDEL\par
NORDISK FORLAG\par\medskip
\rule{1.5cm}{0.4pt}\par\medskip
{\small FORLAGETS BOGTRYKKERI}\par\medskip
1907
\end{titlepage}
% --- p. VII ---
\opage{VII}%

\chapter*{FORORD TIL FEMTE UDGAVE}


Dette Omrids af Logiken følger i Fremstillingsmaade
den af \textsc{Leibniz}, \textsc{Boole} og \textsc{Jevons} opstillede Identitets- og
Indholdslogik. Ved gentagen Syslen med denne Fremstillingsmaade
og Anvendelse af den ved min Undervisning
er jeg bleven bestyrket i, at den baade tager nærmere
Sigte paa den væsentlige Side ved vor Tænken, og at den
% PRINTED AS IS: »pædagisk« (for pædagogisk)
ved sin Simpelhed og Klarhed egner sig bedre til pædagisk
Brug end den sædvanlige Fremstillingsmaade. ---

I nærværende femte Udgave har jeg foretaget en gennemgaaende
Omarbejdelse, støttet til de Synsmaader, jeg
har udviklet i min Afhandling \textit{Det psykologiske Grundlag
for logiske Domme} (Videnskabernes Selskabs Skrifter. 6.
Række) og i Afhandlingen om Begrebet Analogi (Mind
1905, paa Dansk i anden Række af mine „Mindre Arbejder“).
Jeg har tillige kunnet gøre Brug af nogle værdifulde
kritiske Bemærkninger fra en af mine tidligere Tilhørere
og fra Dr. \textsc{Anton Thomsen}.

\textit{December 1906.}

\hfill\textbf{Harald Høffding.}

% not transcribed: the accession stamp "50811" and the dealer's stamp "Dec. 14. 1907. Stechert"

% --- p. IX ---
\chapter*{\opage{IX}INDHOLDSFORTEGNELSE}
\begin{small}


\ix{0em}{\textbf{I. Syntese og Analyse.}}{1}
\par\noindent\textbf{II. Analyse af Anskuelser:} Begrebs- og Domsdannelse.
\ix{1em}{A. Begrebsdannelse}{5}% the layer: "Begrebsdannelse ..." (a stop before the leader not certain)
\ix{1em}{B. Begrebsindhold og Begrebsomfang.}{8}
\ix{1em}{C. Dommes Dannelse og Formulering.}{11}% head p. 11: "C. Dommens Dannelse og Formulering."
\ix{1em}{D. Sekundære Domme.}{18}
\par\noindent\textbf{III. Analyse af Domme:} Slutning.
\ix{1em}{A. De logiske Grundsætninger}{23}
\ix{1em}{B. Slutninger}{28}
\ix{2em}{a. Direkte Slutning}{29}
\ix{2em}{b. Indirekte Slutning.}{34}
\ix{3em}{α. Indirekte Slutning ved ren Deduktion}{35}% head p. 35: "α. Indirekte Slutning som ren Deduktion."
\ix{3em}{β. Indirekte Slutning som Forbindelse af Induktion og Deduktion}{38}
\par\noindent\textbf{Tillæg:}
\ix{1em}{1. Mill’s Induktionsmetoder}{42}% head p. 42: "1. Mill’s induktive Metoder."
\ix{1em}{2. Analogi og Analogislutning}{46}
\end{small}
\clearpage

\mainmatter
% notatex.py applied
% heads.py applied

% --- p. 1 ---
\opage{1}\chaphead{I}{SYNTESE OG ANALYSE}{}


1. At tænke er at sammenligne, at finde Forskel og Lighed.
Under forskellige Former og i forskellige Grader ytrer en
saadan Sammenlignen sig i al Fornemmen og Forestillen, i al
Iagttagelse og Forestillingsforbindelse. Ja, enhver Bevidsthedstilstand
er betinget ved de Forskelle og Ligheder, der findes
indenfor den og i dens Forhold til andre Bevidsthedstilstande.
Lyst og Smerte betinges ved en Oplevelses Forhold til vor Natur
og vor Tilstand; at ville er at foretrække, og beror altsaa ogsaa
paa et Forskelsforhold. I snevrere Betydning kaldes kun den
med udtrykkelig, koncentreret (uvilkaarlig eller vilkaarlig) Opmærksomhed
fuldbyrdede Sammenligning Tænken. En saadan
Opmærksomhed vækkes, naar det gælder at løse en Opgave
eller finde Midler til et Formaal. (Smlgn. Psykologi VB, 11.)

Sammenligning forudsætter Energi til paa én Gang eller i
umiddelbar Rækkefølge at kunne fastholde og sammenholde
forskellige Elementer; kun ved saadan Sammenholden fremtræder
Ligheden eller Forskellen. En psykisk Energi, der ytrer
sig som Syntese, er derfor al Tænknings sidste Forudsætning.
(Smlgn. Psykologi II, 5.)

Ogsaa fra en anden Side set er Syntese Tænkningens sidste
Forudsætning. Sammenligningen sker ved, at Opmærksomheden
gaar fra Element til Element; men der forudsættes da, at de
forskellige Elementer forud for Sammenligningen ere givne i
Sammenhæng, i en Totalanskuelse. Kun paa Grundlag af en
saadan Totalanskuelse kan Analysen, den særlige Henvendelse
af Opmærksomheden paa de enkelte Elementer (Dele, Sider
% --- p. 2 ---
% running head: I. Syntese og Analyse .
\opage{2}eller Beskaffenheder), foregaa. --- Totalanskuelsen kan være
Sanseanskuelse, Erindringsanskuelse eller Fantasianskuelse. Der
kan finde en Tænken Sted paa Grundlag af alle disse forskellige
Arter af Anskuelse, og det bliver da et specielt Spørgsmaal,
hvorledes disse forholde sig til hverandre. --- I enhver
saadan Totalanskuelse er en Mangfoldighed af Elementer forbundne.


Den sammenholdende og forbindende Virken, der saaledes
fra to Sider set er Betingelse for Sammenligning, altsaa for
Tænken, kalde vi Syntese. \textit{Al Analyse forudsætter Syntese,
ligesom al Sammenligning gør det}.

Totalanskuelser dannes fra først af uvilkaarligt, og vi holde
saa længe som muligt fast ved dem. Kun modstræbende skride
vi til at opløse dem. Men en anden Ordning af Elementerne
end den, der er givet i de uvilkaarlige Anskuelser, kan være
nødvendig, dels for at skaffe os en Oversigt over de mange
forskellige Anskuelser, vi have, dels for at ophæve de Modsigelser,
der kunne fremtræde indenfor samme Anskuelse eller
i Forholdet mellem flere forskellige Anskuelser, der optræde
samtidigt eller paa nært Hold.

Medens Psykologien som Tænkningens Naturhistorie undersøger
de faktiske Tankeprocesser, der foregaa i os, søger Logiken
Regler og Maalestokke, efter hvilke der kan skelnes
mellem gyldig og ugyldig Tænken. Den gyldige Tænken maa
være den, der naar til faste og almengyldige Resultater. Ved
et fast Resultat menes et saadant, som kan fastholdes, uden at
den Tænkende kommer i Modsigelse med sig selv (om end
nok saa mange Sanse-, Erindrings- og Fantasianskuelser dukke
op); ved et \textit{almengyldigt} Resultat menes et saadant, som
anerkendes ikke blot af den Tænkende selv, men af andre
tænkende Væsener paa samme Udviklingstrin som hans. Den
logiske Tænken udspringer af en Stræben efter at bevare
Identiteten med os selv og med Alle, hvis Tænkeevne vi
anerkende som lige med vor egen. \textit{Den gyldige Tænken er
den, ved hvilken det er muligt at hævde saa stor Overensstemmelse
som muligt mellem saa mange Bevidsthedselementer
(vore egne og Andres) som muligt}. Erfaringen viser, at jo mere vi
% --- p. 3 ---
% running head: I. Syntese og Analyse .
\opage{3}fastholde den herved givne Maalestok, des færre Skuffelser lide
vi, og des tryggere Forventninger kunne vi hengive os til. ---

Den logiske Tænken er i en vis Forstand ikke produktiv.
Dens Arbejde gaar ud paa gennem Analyse og Sammenligning
at bringe til udtrykkelig Bevidsthed, hvad der allerede er indeholdt
i de givne Totalanskuelser, som ere den oprindelige Synteses
Værk. Den bringer det i vore Sanse-, Erindrings- og
Fantasianskuelser givne Indhold i nye Forbindelser, som kaldes
% PRINTED AS IS: the semicolon before »i hvilke«
Begreber og Domme. Af Domme kan igen drages Slutninger;
i hvilke Forholdet mellem de i forskellige Domme indeholdte
Begreber komme til udtrykkelig Bevidsthed; men et helt nyt
Indhold bringer Slutningen ikke i Forhold til Dommene, saalidt
som Domme og Begreber i Forhold til Anskuelsen. Og
dog kan det ikke blot være et strengt, men ogsaa et betydningsfuldt
psykisk Arbejde, der er gjort, naar Bevidsthedsindholdet
er omsat i nye Former (fra Anskuelser til Begreber og
Domme, og fra Domme til Slutninger), saa at nye Synteser
kunne træde i Stedet for de første Synteser. Derved kan
der ikke blot skaffes Oversigt og Klarhed over det hidtil Anskuede,
men der kan ogsaa vækkes Trang og Evne til at opsøge
og frembringe nye Anskuelser.

2. \textit{Enhver} Videnskab er en Lære om gyldig Tænken paa
sit Omraade. Den formelle Logik undersøger den gyldige
Tænkens almindelige Former, Betingelserne for Tankens Overensstemmelse
med sig selv, ligegyldigt hvilke Omraader Tænkningen
bevæger sig paa. Den ser bort fra alt bestemt Indhold
og kan derfor hente sine Exempler fra alle mulige Omraader.
Det er jo stedse en og samme menneskelig Tanke, der arbejder,
og der maa da være visse Gyldighedsregler, der gælde overalt.

Matematiken kan ogsaa kaldes en formel Videnskab, da den
ser bort fra Genstandenes specielle Kvaliteter og blot undersøger
Ligheder og Forskelle mellem Figurer og Størrelser.
Men Logiken er i endnu højere Grad formel, idet ogsaa Figurer
og Størrelser for den er specielle Kvaliteter; den undersøger
Lighed og Forskel i al Almindelighed, hvorsomhelst de fremtræde
for Tænkningen. De logiske Grundsætninger danne derfor
Forudsætninger for Matematiken som for al anden Videnskab.

% --- p. 4 ---
% running head: I. Syntese og Analyse .
\opage{4}Den formelle Logik falder ikke sammen med Erkendelsesteorien,
som ogsaa undertiden kaldes Logik. Hin prøver Overensstemmelsen
mellem vore Forestillinger uden at spørge om
disses reale Betydning. Logisk Tænkning kan øves paa Grundlag
af Fantasianskuelser saavel som paa Grundlag af Erindrings- og
Sanseanskuelser. Den formelle Logik gaar over til Erkendelsesteorien,
naar dér spørges om, hvorvidt vor Tænknings Former
og Resultater gælde for den reale Tilværelse: om der er Konsekvens
i Tilværelsen, ligesom der er Konsekvens i, hvad Logiken
anser for den rette Tænken. Dette Spørgsmaal kommer
naturligt frem ved, at Tænkningen øves paa Grundlag af Anskuelser,
den ikke selv har frembragt. I Sammenhæng hermed
opstaar ogsaa Spørgsmaalet om Forholdet mellem Sanse-,
Erindrings- og Fantasianskuelser, et Forhold, den formelle Logik
lader staa hen.

% --- p. 5 ---
\opage{5}\chaphead{II}{ANALYSE AF ANSKUELSER}{DANNELSE AF BEGREBER OG DOMME}


\sechead{A. Begrebsdannelse.}{\textbf{A. Begrebsdannelse.}}


3. Den første Betingelse for gyldig Tænken er, at vi blive
os tydeligt bevidste, hvad der indeholdes i de ved uvilkaarlig
Syntese dannede Anskuelser. Dette opnaas ved Analyse, altsaa
ved en Række af Opmærksomhedsakter (se 1). Af det Bord,
jeg sidder ved, har jeg et Billede, der ved en uvilkaarlig
Syntese er dannet af umiddelbare og reproducerede Syns-,
Berørings- og Bevægelsesfornemmelser. Jeg analyserer det, naar
jeg bemærker, at det er sort, haardt, firkantet, langt o. s. v. Et
kemisk Stof (f. Ex. Kulsyre) bestemmer jeg ved at lægge Mærke
til dets Udseende, dets Virkninger, dets Forekomst i Naturen,
dets Opstaaen. I Psykologien karakteriserer jeg Bevidstheden
ved at lægge Mærke til dens Grænsetilstande og til dens ejendommelige
Virkemaade (se Psykologi II, 5).

4. Denne Analyse støttes dels ved Sammenligning mellem
Elementerne (Delene eller Egenskaberne) indbyrdes, dels ved
Sammenligning mellem dem og lignende eller forskellige Elementer
i andre Anskuelser, som reproduceres (dels i bunden,
dels i fri Form, se Psykologi VB, 1---2). Kun derved kan det
enkelte Elements Ejendommelighed ret fremtræde. Den enkelte
Dels Beliggenhed og Størrelse kan kun bestemmes, naar den
sammenholdes med andre Dele; den enkelte Egenskab maa
sammenholdes med andre, forat den kan blive tydelig. (Jvnfr.
Forholdsloven. Psykologi VA, 5; D, 5). Ofte kan et Element
kun skelnes, naar det er forekommet i en anden Sammenhæng,
% --- p. 6 ---
% running head: II . Analyse af Anskuelser .
\opage{6}saa at en Genkendelse kan ske. Bordets Farve og Form lægger
jeg snarere Mærke til, naar jeg har set samme Farve med anden
Form og samme Form med anden Farve. Denne Genkendelse
udtrykkes ved, at jeg giver Egenskaberne Navne. Kalder jeg
Bordet firkantet, vil det sige, at jeg genkender en Form, jeg
har truffet i andre Anskuelsesbilleder. Naar Newton kaldte den
Kraft, der holder Planeterne i Nærheden af Solen, Tiltrækning,
vilde han derved blot betegne, at det er den samme
Kraft, som fører den udslyngede Sten tilbage til Jorden.

Analysen støtter sig saaledes til en Syntese, der virker som
Hjælpeproces. Da allerede Anskuelse og Sammenligning forudsætte
Syntese (1), se vi, hvorledes Analyse og Syntese staa i
stadig Vexelvirkning.

5. Have vi saa fundet de Elementer, der indeholdes i den
givne Anskuelse, gælder det at forbinde og ordne dem saaledes,
at vi faa en samlet Karakteristik af det givne Indhold. Der
maa da igen foretages en Sammenfatten, en Syntese; først
derved afsluttes hele den Proces, der fører til Opstilling af \textit{et
Begreb}.

Fra den blotte Forestilling, der kun er Reproduktion af en
given Anskuelse (i simpleste Tilfælde kun af en Fornemmelse),
adskiller Begrebet sig ved den gennem Vexelvirkningen af
Analyse og Syntese vundne Klarhed og Tydelighed. Et Begreb
er en Forestilling, hvis Indhold vi ere blevne os klart og tydeligt
bevidst, saa at det ikke ændres ved den forskellige Sammenhæng,
i hvilken det forekommer. Vor Bevidsthed kan nu paa
dette Punkt bevare sin Identitet med sig selv, og tilfældige
Forestillingsassociationer kunne hindres fra at gribe ind. Derved
bliver det tillige muligt, at flere Individer kunne tænke i
Forening: kun naar de hvert for sig have fæstnet deres Forestillinger,
kunne de til samme Ord knytte de samme Forestillinger.


\textit{Et Begreb er altsaa en (for samme Individ) fast og (for
de forskellige tænkende Individer) almengyldig Forestilling}.
--- Definitionen (Begrebets Fastsættelse, Begrebsbestemmelsen)
bestaar væsentligt i en Afgrænsning af Forestillingernes Indhold.
Den søger at reducere dette til det mindst mulige ved at ude%
% --- p. 7 ---
% running head: A. Begrebsdannelse .
\opage{7}lade de Elementer, som ikke ere nødvendige til at skelne det
ene Forestillingsindhold fra det andet.

Et Begreb kan være dannet paa Grundlag af en enkelt
Anskuelse, eller paa Grundlag af flere forskellige Anskuelser
af samme Fænomen (Ting, Person, Begivenhed, Forhold), eller
paa Grundlag af flere Anskuelser, der have forskellige Fænomener
til Genstand. I de to første Tilfælde kaldes Begrebet
et \textit{Individualbegreb}, i det sidste Tilfælde et Almenbegreb.
(Smlgn. Psykologi VB, 9.) Dannes Individualbegrebet paa Grundlag
af en enkelt Anskuelse, kaldes det \textit{konkret}; dannes det
paa Grundlag af flere forskellige Anskuelser, kaldes det \textit{typisk}.
Almenbegreber kunne være af lavere eller højere Grad, alt
eftersom de staa umiddelbar Anskuelse nærmere eller fjernere;
hine kaldes \textit{Artsbegreber}, disse \textit{Slægtsbegreber}. Hvad der i en
Henseende er Art, kan i en anden Henseende være Slægt, og
omvendt.

Begrebets \textit{Omfang} er Summen af de Anskuelser, for hvilke
det gælder; Begrebets Indhold er de Kendemærker (Egenskaber
eller Dele), som ere fundne ved Analyse af de givne
Anskuelser.

Hvor flere forskellige Anskuelser ligge til Grund, altsaa
ved typiske og almene Begreber, behøver den Lighed, der
findes mellem de enkelte Anskuelsesbilleder, kun at være Forholdslighed
(Analogi): det vil være Forholdet mellem Delene
eller Egenskaberne, der er ligt, ikke selve de enkelte Dele eller
Egenskaber. Ved Individualbegrebet om et bestemt Menneske,
der skal gælde for det baade som ungt, som voxent og som
gammelt, er Ligheden, hvorpaa Begrebet bygger, Forholdslighed.
Ligeledes ved Begrebet Hest, Begrebet Hvirveldyr o. s. v.
Begrebet udtrykker Loven for Delenes eller Egenskabernes
indbyrdes Forhold. Kun hvor Begrebet angiver et ganske enkelt
Kendemærke (som f. Eks. Mærket paa Kofferter, der indskrives
fra København til London), kan Dækningslighed eller Kvalitetslighed
være tilstede. Hvor dette ikke er Tilfældet, er det kun
ved Analogiens Hjælp, at Bevidstheden kan gaa fra Exempel til
Exempel med Bevarelse af sin Identitet med sig selv. (Smlgn.
om Dæknings-, Kvalitets- og Forholdslighed Psykologi VB, 8.)
% --- p. 8 ---
% running head: II . Analyse af Anskuelser .
\opage{8}Det simpleste og klareste Exempel lægges til Grund, maaske
blot i skematisk Form, d. v. s. med Udeladelse af alle undværlige
Elementer. Begrebet er en Anvisning til at foretage en
Tankebevægelse gennem en Række af Anskuelser. Om den
kan fuldføres, beror paa, om Begrebet er dannet paa Grundlag
af tilstrækkeligt mange og tilstrækkeligt fuldstændige Anskuelser.
Hvor Anskuelsesindholdet er uudtømmeligt, bliver Definitionen
kun en Hypotese, der kan være nødvendig til foreløbig Oversigt,
men maa finde sin Bekræftelse ved fortsat Analyse.

6. Begrebsdannelsen sker ved, at en given Anskuelses Elementer
ses i deres indbyrdes Forhold og i Forhold til Elementerne
i andre Anskuelser, og Begrebsbestemmelsen sker
ved en Afgrænsning af det Indhold, der er fundet ved Analyse.
\textit{Ethvert Begreb udtrykker derfor et Forhold}. Jeg bliver
mig Elementerne i en Anskuelses Indhold bevidst ved at betragte
dem i deres Forskel og Lighed dels overfor hverandre
indbyrdes, dels overfor Elementer i andre Anskuelser. Ethvert
Begreb maa derfor defineres i Forhold til andre Begreber.
En Definition maa ændres, dels naar Anskuelsesindholdet
ændres, dels naar de nærmest liggende Begreber ændres.
Der er derfor en stadig Vexelvirkning mellem Begreb og Anskuelse
og mellem Definition og Klassifikation. Begrebet kemisk
Sammensætning ændredes, da Boyle paaviste, hvorledes der
ved Sammensætning af elementære Stoffer kan opstaa et Stof
med andre Egenskaber end disse; og derved blev det fremdeles
muligt at skelne skarpt mellem kemisk Sammenhæng
og mekanisk Blanding. Rousseau’s Paavisning af Følelseslivets
Ejendommelighed førte til en ny Inddeling af Psykologien.

Kan man ikke paavise Indholdet af et Begreb, der danner
Grænse for et givet Begreb, danner man foreløbigt et negativt
Begreb (non $A$ som Grænse for $A$). Se nedenfor: 8.

\sechead{B. Begrebsindhold og Begrebsomfang.}{\textbf{B. Begrebsindhold og Begrebsomfang.}}


7. \textit{Begrebers Indhold og Omfang} (5) \textit{staa i omvendt Forhold
til hinanden}. Ti jo flere forskellige Anskuelser, der ligge
til Grund ved Begrebsdannelsen, des færre Kendemærker kunne
% --- p. 9 ---
% running head: B. Begrebsindhold og Begrebsomfang .
\opage{9}optages i Begrebets Indhold. Jo flere forskellige Anskuelser
et Begreb skal gælde for, des færre Egenskaber kan det udsige
om dem. Saaledes har det typiske Individualbegreb større
Omfang, men mindre Indhold end et tilsvarende konkret Individualbegreb;
et Almenbegreb har større Omfang, men mindre
Indhold end et Individualbegreb, og paa samme Maade gaar det
med Arts- og Slægtsbegreber i Forhold til hinanden. Et Begreb
kaldes lavere i Forhold til et andet, naar det har mindre Omfang,
men større Indhold, --- højere, naar det har større Omfang,
men mindre Indhold. Saaledes er Begrebet Menneske
højere i Forhold til Begrebet Neger: det omfatter flere Individer
end dette, men det kan ikke angive saa mange Kendemærker
som dette.

Sætningen om det omvendte Forhold mellem Omfang og
Indhold gælder kun, naar vi sammenligne Begreber, der ligge
i samme Række, d. v. s. naar de finde Anvendelse indenfor en
og samme Gruppe af Anskuelser og Kendemærker, uden dog
at være sideordnede. Den gælder ikke for to sideordnede Artsbegreber;
ti en Arts Omfang kan udvides eller indskrænkes,
uden at dens Indhold forminskes eller forøges. Fordi der er
flere Spurve end Ørne i Verden, behøver Begrebet Spurv
ikke at have mindre Indhold end Begrebet Ørn. Begrebet
Fugl har derimod større Omfang, men mindre Indhold end
begge de to nævnte Begreber. Andre Exempler paa Begreber,
der ligge i samme Række, ere: Neger --- Menneske --- levende
Væsen; Had --- Følelse --- Bevidsthedselement; Kvadrat --- Rektangel
--- Firkant --- retliniet Figur --- Figur.

Det maa tillige bemærkes, at vi ikke altid tænke et Begreb
med alle dets Kendemærker, men kun med dem, der ere nødvendige
for at skelne det fra de nærmest liggende Begreber.
Vi lade os nøje med at definere Mennesket som Fornuftvæsen,
naar det blot gælder om at skelne Mennesket fra Dyret. Loven
om det omvendte Forhold mellem Indhold og Omfang gælder
naturligvis kun, naar vi tage Hensyn til hele Indholdet.

8. \textit{Negative Begreber} dannes, naar det ved Forsøg paa
Definition viser sig, at der er en Grænse for et Begrebs Indhold
og Omfang, uden at man foreløbigt kan angive positive
% --- p. 10 ---
% running head: II . Analyse af Anskuelser .
\opage{10}Kendemærker for det, der ligger paa den anden Side af Grænsen.
Ved Klassifikationen maa vi da lade os nøje med at opføre
et Begreb med negtende Kendemærker: non $A$ som Grænse
for $A$. (Denne Betegnelse minder om de i 5 omtalte Begreber,
hvis eneste Indhold er et simpelt Mærke.) Klassifikationernes
Historie viser, at man ofte foreløbigt maa danne Grupper uden
positive Kendemærker og overlade til den senere Forsken at
finde saadanne, hvilket da vil føre til en videre Inddeling af
den Mangfoldighed, der maaske indbefattes i det negative Begreb.
Aristoteles inddelte Dyrene i fornuftige og ufornuftige,
og de ufornuftige igen i dem, der have (rødt) Blod, og dem, der
ikke have (rødt) Blod. Ofte kan et negativt Begreb skjules
under en positiv Benævnelse; saaledes kunde Linné for sin
sjette Dyreklasse, som han kaldte Ormene, ikke anføre andre
Kendemærker, end at de hverken vare Pattedyr, Fugle, Amfibier,
Fisk eller Insekter. Senere inddelte Lamarck Dyrene i
Hvirveldyr og Hvirvelløse, og først Cuvier fandt positive
Kendemærker for og en positiv Inddeling af de Hvirvelløse.
--- Negative Begreber kunne altsaa have den Betydning at markere
foreløbige Grænser under Forskningens Udvikling.

Forholdet mellem to negative Begrebers Indhold er det
omvendte af Forholdet mellem de to tilsvarende positive Begrebers
Indhold. Begrebet „levende Væsen“ har mindre Indhold
end Begrebet „Menneske“, men Begrebet „Ikke-levende“ har
større Indhold end Begrebet „Ikke-Menneske“. Altsaa $A < B$,
men $a > b$ (naar de smaa Bogstaver betegne de negative Begreber).
--- Det Tilsvarende gælder for Omfanget. Begrebet
„levende Væsen“ har større Omfang end Begrebet „Menneske“,
men Begrebet „Ikke-levende“ har mindre Omfang end Begrebet
„Ikke-Menneske“. $A > B$, men $a < b$. --- De højere Begrebers
Negationer have større Indhold, men mindre Omfang end de
lavere Begrebers Negationer.

9. Ved Brugen af Ordene „højere“ og „lavere“ er der taget
Hensyn til Begrebernes Omfang, og den ældre Logik har
idethele lagt overvejende Vægt paa denne Side af Begrebet.
Og dog er det Indholdet, der først og især spiller en Rolle
% --- p. 11 ---
% running head: B. Begrebsindhold og Begrebsomfang .
\opage{11}ved Tænkningen. Det maa først afgøres, hvilke Kendemærker
der skulle optages i Begrebet, før det kan afgøres, i hvor
mange forskellige Tilfælde disse Kendemærker genfindes.
Dette hænger sammen med, at vi fra først af danne vore Begreber
paa Grundlag af Anskuelse; naar vi gaa strengt metodisk
frem, maa vi derfor begynde med konkrete Individualbegreber.
Det staar fra først af ubestemt, hvor stort Omfanget
kan blive, --- om vi kunne danne typiske og almene Begreber.
Det er altsaa Indholdet, vi begynde med, naar vi gaa logisk
frem. --- Psykologien lærer os ganske vist, at Kunsten netop
kan være at opfatte det konkret Individuelle i dets hele Ejendommelighed,
og at en vag Opfattelse af det Enkelte forleder
til forhastet Generalisation, saa at den Udvikling, vore Forestillinger
maa gennemgaa, fra en væsentlig Side set bestaar i
en Individualisering og Specialisering. (Psykologi VB, 9b.) Men
det viser kun, at den logiske Lære om den gyldige Tænken er
forskellig fra den psykologiske Lære om den faktiske Tænken.

\sechead{C. Dommens Dannelse og Formulering.}{\textbf{C. Dommens Dannelse og Formulering.}}


10. Domsdannelse og Begrebsdannelse kunne ikke skilles
fra hinanden, men ere to Sider af samme Tankeproces. Ogsaa
Dommen dannes, hvor den dannes fra først af, ved Analyse
af Anskuelse. Medens Analysen ved Begrebsdannelsen gaar ud
paa Fremhæven og Udskillen af enkelte Elementer, der komme
til at udgøre Begrebets Indhold, gaar den ved Domsdannelsen
ud paa at bestemme Forholdet mellem de i Anskuelsen givne
Elementer. Men til et tydeligt Begreb hører ogsaa Bestemmelsen
af dette Forhold, og selve Opstillingen, Definitionen
af Begrebet, sker ved en Dom, idet det enkelte Begreb kun
kan bestemmes ved sit Forhold til andre Begreber. (Smlgn. Forholdet
mellem Definition og Klassifikation: 6. 8.) Den hjælpende
Syntese (4) spiller en større Rolle ved Domsdannelsen end
ved Begrebsdannelsen, da Opfattelsen af det indbyrdes Forhold
mellem Anskuelsens Elementer saavel som af Anskuelsens Forhold
til andre Anskuelser forudsætter mere udpræget Sammenholden
med Noget, som ikke umiddelbart er givet.

% --- p. 12 ---
% running head: II . Analyse af Anskuelser .
\opage{12}En Dom er en \textit{fast og almengyldig Forbindelse af Begreber}.
(Smlgn. 5.)

Efter Sagnet skal Newton ved Synet af et faldende Æble
være kommen paa den Tanke, at Tyngden maatte virke i hele
Verdensrummet, altsaa paa Maanen i Forhold til Jorden og paa
Planeterne i Forhold til Solen, og af det saaledes opstaaede
Anskuelsesbillede har han uddraget den Dom, som han senere
søgte at bevise: Planetbevægelsen er en Faldbevægelse. --- Siger
jeg: denne Mand er god! --- udledes denne Dom (naar jeg
fælder den paa første Haand) af mine Sanse- og Erindringsanskuelser
af Manden og hans Færd, i hvilken jeg finder Egenskaben
„god“ (som jeg kender fra andre Anskuelser, se 4).

Ved et Bevis for Dommens Gyldighed maa man gaa Slutningens
Vej, gennem Induktioner og Deduktioner. Forholdet
mellem Anskuelse og Dom er ganske vist analogt med Forholdet
mellem Grund og Følge, men er ikke saa fuldt og klart
bevidst, som dette kan være.

Denne Analogi finder ogsaa Sted i den Henseende, at ligesom
Følgen kun indeholder, hvad der allerede indeholdtes i
Grunden, saaledes bringer Dommen ikke noget helt nyt Indhold,
men kun en ny, mere bestemt og bevidst Form for det
i Anskuelsen givne Indhold.

11. Det Forhold, Dommen udsiger som gældende mellem
to Begreber (eller Begrebselementer), kan udtrykkes paa forskellig
Maade. Man kan holde sig til \textit{den sædvanlige sproglige
Udtryksmaade}, hvor det ene Begreb eller Element (Prædikatet)
henføres til det andet (Subjektet) ved Hjælp af Ordet
„er“ (Kopula). Eller man kan holde sig til, at det ene Begreb
(Prædikatsbegrebet) har et større Omfang end det andet (Subjektsbegrebet),
saaledes at dette indbefattes i (subsumeres under)
hint. Dette \textit{Subsumptionsforhold} kan man da anskueliggøre
ved geometriske Figurer, saa at f. Ex. Subjektet fremstilles som
en mindre Cirkel inde i en større, der betegner Prædikatet.
(En analog Fremstillingsmaade kunde anvendes med Hensyn til
Forholdet mellem Begrebernes Indhold, saa at Prædikatsbegrebet
udtryktes som en lille Cirkel inde i Subjektsbegrebets større
Cirkel. Denne Fremstillingsmaade er dog aldrig bleven forsøgt.)
% --- p. 13 ---
% running head: C. Dommens Dannelse og Formulering .
\opage{13}Eller man kan formulere Dommen som Udtryk for \textit{et Identitetsforhold}
mellem to Begreber, altsaa som en Ligning. Denne
Udtryksmaade minder om Algebra, ligesom den forrige om
Geometri.

Men skønt Dommen kan minde baade om Grammatik,
Geometri og Algebra, er den dog i sig selv uafhængig af disse
Udtryksmaader, hvad der ses af, at Ord, Figurer eller Tegn
maa vælges efter Tankeforholdet og bestemmes ved dette. En
nærmere Paavisning heraf vil tjene til Belysning af Tankens
Forhold til den Maade, den udtrykkes paa.

12. Det logiske Subjekt er det mere ubestemte Begreb,
hvorfra Tankebevægelsen ved Domsdannelsen gaar ud, og
hvorom den drejer sig (Udgangsforestillingen); det logiske Prædikat
er det Begreb, i hvilket Tankebevægelsen ender (Slutningsforestillingen),
fordi Subjektet i det har fundet sin nærmere
Bestemmelse. Det danske Sprog forbinder dem ved Ordet „er“.
Men den grammatiske Sætning er ikke altid en simpel Formulering
af et saadant Begrebsforhold.

a) \textit{En grammatisk Sætning kan paa én Gang indeholde
flere logiske Domme}. Sætningen „Denne Sofa er haard“ kan
indeholde tre Domme, idet Udsagnet kan stille denne Sofa i
Modsætning til en anden Sofa, eller betegne den paapegede
Genstand som en Sofa, eller blot udsige Haardheden om den.
„Gustav Adolf faldt ved Lützen den 6. November 1632“ udsiger
ikke blot Kongens Fald, men Tid og Sted. b) Denne
Flerhed af Domme i en og samme Sætning hænger sammen
med, at \textit{det grammatiske Prædikat ikke altid er det logiske
Prædikat}. Man vil kunne finde det logiske Prædikat ved at
bemærke, hvor Tonen ligger. Siges „\textit{denne} Sofa er haard“,
ligger det logiske Prædikat i „denne“, og Dommen bliver: „den
Sofa, som er haard, er denne her.“ „\textit{Alle} Gæsterne ere komne“:
her ligger det logiske Prædikat i „Alle“; det er det Nye. „\textit{Da}
sprang Krudttaarnet“: det logiske Prædikat er her Tidspunktet:
man vidste i Forvejen, at Krudttaarnet var sprunget i Luften.
Ved Sætninger, der ikke udtales mundtligt, og hvor Betoningen
ikke angives, maa Fortolkning ud af hele Sammenhængen finde
det logiske Prædikat. c) \textit{Der gives Sætninger, i hvilke det}
% --- p. 14 ---
% running head: II . Analyse af Anskuelser .
% PRINTED AS IS: »af« printed twice (line end and the next line: »af / af sig selv«)
\opage{14}\textit{logiske Subjekt slet ikke er udtrykt}, fordi det forstaar sig af
af sig selv eller ligger i Situationen. Saadanne Prædikatsdomme
have vi i de simple Udraab, som „Herligt!“ „Bravo!“, eller i
Meldinger som „Kongen!“ --- Herhen høre ogsaa de ubestemt
beskrivende Sætninger, i hvilke blot en Tilstand eller en Virksomhed
angives: „Det trækker.“ „Det gaar tilbage.“ „Der tales
om ....“ Udgangsforestillingen er her ubestemt og kan ikke
formuleres, eller den er selvfølgelig, og Interessen gaar ganske
op i det logiske Prædikat.

Betoningens Betydning og Tilstedeværelsen af Prædikatsdomme
viser, at den psykologiske Interesse er knyttet til det
logiske Prædikat, som bringer det Nye, medens det logiske
Subjekt staar som foreløbigt bekendt, selvom det trænger til
nærmere Bestemmelse. Skal der naas en klar Bevidsthed om
det Forhold, Dommen hævder, maa denne dog formuleres saaledes,
at baade Subjekt og Prædikat bestemmes udtrykkeligt;
kun da kan Forholdet mellem dem blive klart, det Forhold,
hvortil den logiske Interesse egentligt er knyttet. Kun naar
dette Forhold er tydeligt angivet, kan Dommen virkeligt være
en fast og almengyldig Forbindelse af Begreber. Logisk set er
det derfor tilsidst uden Betydning, hvad der er Udgangsforestilling,
og hvad der er Slutningsforestilling, og hvorledes vi
komme fra hin til denne. Tidsforholdet taber sin Betydning.
Den grammatiske Formulering staar i denne Henseende nærmere
ved Psykologien end ved Logiken. De to følgende Formuleringer
have derimod en rent logisk Karakter.

13. Naar man, som den ældre Logik ligefra Aristoteles af,
lægger Begrebernes Omfang til Grund, betyder Dommen, at
den Gruppe af Fænomener, som det ene Begreb (Subjektsbegrebet)
gælder for, er en Del af den Gruppe af Fænomener,
det andet Begreb (Prædikatsbegrebet) gælder for. Dommen udtrykker
altsaa \textit{et Underordningsforhold (Subsumptionsforhold)}.
At Negrene ere Mennesker vil sige, at Gruppen Negre er mindre
end Gruppen Mennesker og indbefattes i denne. Man kan anskueliggøre
dette Forhold ved at fremstille det højere Begreb som
en større Cirkel, indenfor hvis Omkreds det andet Begreb fremstilles
som en mindre Cirkel. At $A$ er $B$, betyder altsaa $A < B$.

% --- p. 15 ---
% running head: C. Dommens Dannelse og Formulering .
\opage{15}Ad denne Vej vil man kunne give en gyldig Fremstilling
af Logiken; kun saadanne Domme, i hvilke de to Begreber
ere absolut identiske ($A = B$), vil man ikke kunne finde Plads
til, og de forekomme heller ikke hos Aristoteles. Geometrisk
vilde vi her kun have en eneste Cirkel og ikke kunne se, at
den betyder to Begrebers Dækning. ---

Fordi man gaar ud fra Begrebernes Omfang, behøver man dog
ikke at opfatte Dommen som Subsumption. I $A < B$ ligger nemlig,
at $A$ er identisk med en Del af $B$’s Omfang. Negrene er en Del
af Menneskeheden. Lad nu $n$ betyde „nogle“, saa er $A = nB$. For
$n = 1$ (hvor 1 betyder det hele Indbegreb i Modsætning til en
Del eller nogle) faas $A = B$. --- Herved fremstilles Dommen
altsaa som Udtryk for et Identitetsforhold mellem to Omfang.

Men det ligger da nær at gaa et Skridt videre. Naar $A$ er
$B$, er $A$’s Omfang maaske kun en Del af $B$’s Omfang; men
i hvert Tilfælde maa $B$’s hele Indhold gælde for $A$. Og da det
nu er Begrebsindholdet, Tænkningen først vender sig imod,
medens Omfangets Bestemmelse er sekundær (9), vil det være
naturligt at lægge Begrebernes Indhold til Grund. Dette fører
os til den tredie Fremstillingsmaade, der udtrykker Dommen
ved et Identitetsforhold mellem Begrebsindhold.

14. Ligesom Trangen til at danne Begreber udspringer af
Bevidsthedens Trang til at bevare sin Overensstemmelse, om
muligt sin Identitet med sig selv trods Bevidsthedsindholdets
Mangfoldighed og Skiften, saaledes vil der ogsaa være en naturlig
Trang til at bevare Identiteten, naar vi forbinde et Begreb med
et andet i en Dom. Jo mere de to Begreber staa udvortes og
fremmede for hinanden, des mindre vil Bevidstheden kunne
bevare sin Enhed, --- des mere vil Forbindelsen staa som paatvungen,
tilsidst som en blot Ordforbindelse. Jo mere vi derimod
kunne se Forholdet mellem de to Begreber som et
Identitetsforhold, des mere naturlig og klar vil Forbindelsen
være. Naar vi nærmere bestemme de forskellige Grader eller
Arter af Identitet, som kunne være tilstede, ville vi kunne
hævde den Paastand, at \textit{Dommen udtrykker et Identitetsforhold}.
--- Leibniz og (i det 19. Aarhundrede) Boole og Jevons have
grundlagt denne Fremstillingsmaade.

% --- p. 16 ---
% running head: II . Analyse af Anskuelser .
\opage{16}Der gives tre forskellige Identitetsforhold.

a) Prædikatsbegrebets Indhold kan være en nærmere Bestemmelse
af Subjektsbegrebets Indhold, uden at de to Begreber
falde ganske sammen. Naar jeg f. Ex. ved at betragte en Mands
Adfærd føres til at udtale den Dom „Manden er god“, er Begrebet
„god“ en nærmere Bestemmelse ved Individualbegrebet
„Manden“. Godheden findes i Mandens Adfærd, og kun fordi
den er der, udsiges Dommen. I Benævnelsen „god“ ligger et
Henblik (ved hjælpende Syntese se 4 og 10) til andre Fænomener,
som have frembudt en lignende Egenskab. Ved Dommen
„Marsbanen er en Ellipse“ udtalte Kepler, efter en Analyse
af Tycho Brahes Iagttagelser, en bestemtere Karakteristik
af Planetens Bane, end man hidtil havde havt; ved Benævnelsen
„Ellipse“ udtryktes en Genkendelse (formedelst hjælpende
Syntese) af en Figur, der var bekendt fra Læren om
Keglesnittene. --- Denne Art Domme kan formuleres saaledes:
$A = AB$. Lighedstegnet betyder her Mulighed af Substitution,
d. v. s. af at sætte det ene Udtryk i Stedet for det andet. Sammenstillingen
af Bogstaverne ($AB$) betyder et Begreb, hvis
Indhold er Summen af to Begrebers ($A$’s og $B$’s) Indhold,
altsaa en Begrebskombination. $A = AB$ betyder altsaa, at naar
vi have $A$, have vi et rigere Indhold, end vi tænkte. $A$ forekommer
ikke uden $B$, men $B$ kan forekomme uden $A$. Der
er andre Ellipser foruden Marsbanen, og andre gode Væsener
foruden hin Mand. $B$ har et større Omfang end $A$: $A < B$.
--- Maaske har $A$ flere Bestemmelser end $B$, og vi have da
$A = ABCD$ .....

Et Identitetsforhold af denne Art kaldes \textit{partiel Identitet}.

b) Naar ikke blot $B$ er en nærmere Bestemmelse ved $A$,
men $A$ ogsaa viser sig at være en nærmere Bestemmelse ved
$B$, saa at vi have baade $A = AB$ og $B = BA$, er der \textit{absolut
Identitet} mellem $A$ og $B$. Det er ligegyldigt, i hvilken Orden
Begreberne forekomme i en Kombination, da dette ikke forandrer
Summen af Begrebsindholdet. Derfor er $AB = BA$.

To absolut identiske Begreber ere kun lige gyldige Udtryk
for et og samme Begrebsindhold, som kan betragtes fra to
forskellige Sider. Exempler: „Dette Hus er mit Hjem.“ „Det
% --- p. 17 ---
% running head: C. Dommens Dannelse og Formulering .
\opage{17}er Peter, der kommer“ (den Kommende = Peter). Kongen af
England = Kejseren af Indien. Grundtvigs Fødselsaar = det
fjerde Aar efter Oehlenschlägers Fødselsaar. Cirkelen = den
i sig selv tilbageløbende Linie, i hvilken ethvert Punkt har
samme Afstand fra et og samme Punkt.

Medens Paavisningen af partiel Identitet kun forudsætter
Analyse af det ene Begreb ($A$), idet det andet ($B$) indeholdes
i det, forudsætter Paavisningen af absolut Identitet en Analyse
af begge Begreber. Den Tankevirksomhed, ved hvilken absolut
Identitet opdages, er altsaa mere sammensat end den, ved
hvilken partiel Identitet opdages. En absolut Identitet forudsætter
to partielle Identiteter ($A = AB$ og $B = BA$). Alligevel
er der en uvilkaarlig Tilbøjelighed til at opfatte et Identitetsforhold
som absolut, før det er prøvet fra begge Sider.

c) Forholdet mellem to Begreber ($A$ og $B$) kan endeligt
være det, at $A$ kun under visse Betingelser er partielt eller
absolut identisk med $B$. Sætningen „Vand er flydende“, gælder
kun, naar Varmegraden er over 0°, men ikke over 80° (Réaumur).
--- Dette Forhold kaldes \textit{begrænset Identitet}. Det udtrykkes
saaledes: $AC = ABC$ (begrænset partiel Identitet) eller
$AC = BC$ (begrænset absolut Identitet). $C$ betegner her den
Betingelse, uden hvilken Identitetsforholdet ikke gælder.

Ved nærmere Undersøgelse vil det vise sig, at enhver Dom
udtrykker begrænset Identitet (ligesom ethvert Begreb udtrykker
et Forhold, se 6). Vi udelade Begrænsningen, hvor den ikke
træder tydeligt frem, og hvor den ikke er af væsentlig Betydning
for Tankegangen. Ofte behøves særegne Erfaringer, forat
Begrænsningen skal opdages. En Tropebeboer vilde ikke mene,
at Sætningen „Vand er flydende“ behøvede nogen Begrænsning
(naar han ikke kalder Damp for Vand).

De Domme, der udsige begrænset Identitet, forudsætte ikke
blot en Analyse af et eller begge Begreber, men ogsaa af hele
% PRINTED AS IS: »deres Indhold forekomme« (singular noun, plural verb)
den større Sammenhæng, indenfor hvilken deres Indhold forekomme,
og indenfor hvilken ogsaa Betingelserne for deres
Identitet maa findes.

% --- p. 18 ---
% running head: II . Analyse af Anskuelser .
\opage{18}\sechead{D. Sekundære Domme.}{\textbf{D. Sekundære Domme.}}


15. Naar der paa Grundlag af Anskuelse er dannet en
Dom, kan der spørges om Gyldigheden af denne Dom, --- om
den gælder for alle Fænomener af samme Art, --- om den
gælder ubetinget, --- om den overhovedet gælder, --- om den
ikke skulde kunne drages i Tvivl. Naar en Dom bestemmes
nærmere i en af disse Henseender, faa vi en Dom om den,
altsaa en Dom i anden Potens. En saadan Dom om en Dom
kaldes en \textit{sekundær Dom}, og i Modsætning til den kaldes den
første Dom en primær Dom. Man har opstillet forskellige
Klasser af sekundære Domme og været tilbøjelig til at betragte
dem som aldeles selvstændige Arter af Domme, der ikke skulde
kunne reduceres til hverandre. Men det vil vise sig, at deres
Forskelligheder for en stor Del kun bero paa de sproglige
Udtryk, og at de, naar vi holde os strengt til deres logiske
Betydning, kunne formuleres som primære Domme. Det er
samme Art Domme, vi have for os, hvad enten deres Subjekt
er et enkelt Begreb eller en Dom.

16. Da Anskuelse (Sansning, Erindring, Fantasi) stedse giver
os noget Enkelt og Individuelt, vil en primær Dom kun angaa
et enkelt Tilfælde. Tankevirksomheden opdager først et partielt,
absolut eller begrænset Identitetsforhold mellem to Begrebers
Indhold; senere maa det afgøres, om dette Forhold
kan antages at gælde i flere Tilfælde, eller maaske i alle Tilfælde
af lignende Art. I $A = AB$ er $A$ et Begreb, hvis Omfang
er mindre end $B$’s (smlgn. 14 a). $A$ er mere individuelt end $B$.
Undersøgelsen af, om $A$ kan gøres til et typisk eller alment
Begreb, sker gennem en Række af Domme, i hvilke Subjektsbegrebets
Indhold sammenlignes med andre, beslægtede Begrebers.
Jeg maa f. Ex. undersøge, om andre Mennesker vise
samme Adfærd som den Mand, jeg (i 14 a) kaldte god, saa at
jeg kan opstille den Dom: „nogle Mennesker ere gode,“ eller
maaske „alle Mennesker ere gode.“ Foretages denne Udvidelse
af Subjektsbegrebets Omfang, bliver den Kvantitetsbestemmelse
(nogle eller alle), der indføres, et logisk Prædikat ved Dommen,
hvorledes den saa end udtrykkes grammatisk. En \textit{kvanti%
% --- p. 19 ---
% running head: D. Sekundære Domme .
\opage{19}tativ} Dom er en Dom om, i hvor mange Tilfælde en vis given
Dom gælder; i denne sekundære Dom er den givne Dom det
logiske Subjekt. Grammatisk føjes Kvantitetsbestemmelsen til
den givne Doms sproglige Subjekt (smlgn. 12: „Alle Gæsterne
ere komne“, ɔ: de nu ankomne Gæster udgøre alle dem, der
ventedes). --- Ofte maa dog Subjektsbegrebets Indhold forminskes,
før Dommen kan udvides til at gælde for flere Fænomener:
dette bestemte Menneskes Begreb har et større
Indhold end Begreberne „nogle“ og „alle“ Mennesker.

En kvantitativ Dom kan udtrykkes ved begrænset Identitet.
Har jeg fundet, at Dommen $A = AB$ kun gælder i nogle Tilfælde,
kan jeg betegne dette ved $CA = CAB$, idet $C$ betyder de
for disse Tilfælde ejendommelige Egenskaber.

Der er ingen Grund til at opstille kvantitative Domme som
en særlig Art af Domme.

17. Naar en Dom fremsættes som ubetinget gyldig, kaldes
den \textit{kategorisk}; naar den fremsættes som betinget ved en
anden Doms Gyldighed, kaldes den \textit{hypotetisk}.

Ved nærmere Undersøgelse viser det sig, at enhver Dom
kan formuleres som en hypotetisk Dom, idet der stedse vil
kunne paavises Forudsætninger, som maa være tilstede, forat
Dommen skal gælde. (Smlgn. 14 c.) Hvis man vil formulere
disse Forudsætninger, kan det ske ved at opstille et begrænset
Identitetsforhold mellem de i Dommen forbundne Begreber:
i $AC = BC$ betegner $C$ Indbegrebet af de Forudsætninger, som
maa være tilstede, forat $A = B$ skal gælde. Dog kan en hypotetisk
Dom ogsaa udtrykkes paa anden Maade. Den hypotetiske
Dom „Hvis $A$ er $B$, er $C D$“ udsiger, at Dommen „$C$ er
$D$“ kun gælder, hvis Dommen „$A$ er $B$“ gælder. Den er altsaa
en Dom om en Dom, en sekundær Dom. De to Domme i en
hypotetisk Dom forholde sig som Subjekt og Prædikat i en primær
Dom. Betegne vi nu Betingelserne for, at $A$ er $B$, ved $P$ og
Betingelserne for, at $C$ er $D$, ved $Q$, saa udsiger hin hypotetiske
Dom en Identitet af $P$ og $Q$: enten $P = PQ$ eller $P = Q$ eller
$PR = PQR$ (eller $QR$), idet $R$ i sidste Tilfælde betegner de
længere tilbage liggende Betingelser for en Identitet af $P$ og $Q$.

Ofte er Forskellen mellem en kategorisk og en hypotetisk
% --- p. 20 ---
% running head: II . Analyse af Anskuelser .
\opage{20}Dom rent sproglig. Sætningen: „Naar Siderne i et Triangel
ere ligestore, ere Vinklerne det ogsaa,“ kan ogsaa udtrykkes
saaledes: „Et ligesidet Triangel er ligevinklet.“

18. De umiddelbare Domme ere positive. Der hører større
Erfaring (flere Anskuelser) til at opstille \textit{negtende} Domme end
til at opstille positive Domme. En negtende Dom opstilles kun,
naar der er Anledning til at afvise en positiv Forventning.
Naar jeg kommer ind i en Stue og siger: „Her er Ingen!“ er
denne negtende Dom en Afvisning af en Forventning om, at
der var Nogen derinde, og den skyldes en Sammenligning
mellem Stuen, som den frembyder sig for mit Blik, og Stuen,
som jeg ventede at finde den. Der viser sig at være færre
Elementer i Iagttagelsen end i Forventningen. (Smlgn. om
negtende Begreber: 8.) En negtende Dom er altsaa en Dom
om en Dom, en sekundær Dom.

En negtende Dom (f. Ex. $A$ er ikke $B$) kan ved nærmere
Betragtning formuleres paa en ligesaa simpel Maade som en
positiv Dom. Den udtrykker nemlig et Identitetsforhold (partielt
eller absolut eller begrænset) mellem det ene Begreb og
det andet Begrebs Negation. Lade vi Tegnet $\sim$ betyde „er ikke“,
og betegne vi (ligesom i 8) ved smaa Bogstaver Negationerne
af de ved de tilsvarende store Bogstaver betegnede Begreber
(saa at f. Ex. $a$ betyder Fraværelsen af de ved $A$ betegnede
Kendemærker), saa er, hvis $A \sim B$, $A = Ab$ og $B = Ba$. Der
er altsaa partiel Identitet mellem $A$ og $b$, og mellem $B$ og $a$.
$A$ og $B$ siges da at staa i \textit{kontrært Forhold} til hinanden: de
ere uforenelige, de udelukke hinanden; – men fordi det ene
Begreb ikke gælder, behøver det andet Begreb ikke at gælde.
Der er flere end to Muligheder. Hvad der ikke er rødt, behøver
derfor ikke at være blaat. Hvad der ikke er godt, behøver
derfor ikke at være ondt. Et retvinklet Triangel er ikke
ligesidet; men fordi et Triangel ikke er ligesidet, er det ikke
sagt, at det er retvinklet --- $A \sim B$ kan dog undertiden betyde
$A = b$. Exempler derpaa haves overalt, hvor der maa træffes
Valg mellem to Muligheder. Efter den strenge stoiske Opfattelse
var Alt ondt, hvad der ikke var godt. Forholdet mellem Position
og Negation (Ja og Nej) er af denne Art; et Spørgs%
% --- p. 21 ---
% running head: D. Sekundære Domme .
\opage{21}maal maa (naar man anerkender de Forudsætninger, under
hvilke det stilles) besvares enten bekræftende eller benegtende.
Negerer jeg $A$, faar jeg $a$; negerer jeg $a$, faar jeg $A$ igen; jeg
kommer altsaa ikke ud over disse to Muligheder. Et saadant
Forhold kaldes \textit{kontradiktorisk}. Der er da ikke blot et Modsætningsforhold,
men et Modsigelsesforhold mellem de to Begreber.
Ikke blot udelukke de hinanden, men naar det ene
ikke gælder, gælder nødvendigvis det andet. Et kontradiktorisk
Forhold er et toleddet Forhold, medens et kontrært Forhold
er et flerleddet Forhold.

Undertiden kan et Forhold, der under sædvanlige Omstændigheder
er kontrært, under visse særlige Betingelser gaa
over til at blive kontradiktorisk. Vi faa da $AC = bC$. Man
maa give Stoikerne Ret i, at man under visse Betingelser
handler ondt, hvis man ikke handler godt, f. Ex. hvis man kan
redde et Menneske, og ikke gør det.

19. Naar der overfor en given Dom kan stilles en anden
Dom, hvis Gyldighed udelukker hins Gyldighed, idet der er et
kontrært eller kontradiktorisk Forhold mellem de to Domme,
--- saa bliver den givne Dom \textit{problematisk} og gør selv den
anden Dom problematisk. De to Domme, Galilei drøfter i sine
„Dialoger“: „Jorden bevæger sig“ – „Jorden bevæger sig ikke“,
gjorde hinanden problematiske. I dette Exempel er Forholdet
mellem de to Domme kontradiktorisk. Derimod have vi et
kontrært Forhold, naar de to Domme: „Bordet er sort“ og
„Bordet er brunt“ stilles imod hinanden; ti her er der flere
end to Muligheder.

Problematiske Domme ere sekundære. De ere egentligt en
Art hypotetiske Domme. En Dom er problematisk, naar dens
Gyldighed forudsætter en anden Doms Ugyldighed, og den kan
derfor formuleres ligesom en hypotetisk Dom med negativ Forudsætning
(altsaa $P = Pq$ o. s. v., se 17). Dog staa ved en problematisk
Dom de to Domme mere Side om Side end ved den
hypotetiske, hvor Forholdet ikke altid er gensidigt (uden ved
$P = Q$). Et Problem opstaar, naar to hinanden udelukkende 
Domme ere mulige, saa at man staar overfor et Enten---Eller.

Kan der føres Bevis for en af Dommene, falder den anden
% --- p. 22 ---
% running head: II . Analyse af Anskuelser .
\opage{22}bort, og den problematiske Karakter ophæves. Det Samme kan
opnaas, naar en af Dommene gendrives. Men hvis Forholdet
er kontrært (ikke kontradiktorisk), kan der være flere Domme,
som stride imod den givne, og som derfor ogsaa maa skaffes
af Vejen, hvis denne skal miste sin problematiske Karakter.
--- En Dom, der paa en af disse Maader har mistet sin problematiske
Karakter, kaldes apodiktisk.

20. Fælles for alle sekundære Domme (de kvantitative,
hypotetiske, negative og problematiske) er det, at den hjælpende
Syntese, der ved en primær Dom kun er af underordnet Betydning
(10 og 14 a), spiller en mere fremtrædende Rolle.
Naar jeg skal bestemme det Omfang, i hvilket en Dom gælder
(16), og de Forudsætninger, under hvilke den gælder (17), maa
jeg udtrykkeligt sammenholde og sammenligne den med andre
Anskuelser eller Domme. Og det Samme er nødvendigt, naar en
Dom skal ophæves (18), og naar dens Modsætningsforhold til
en anden Dom skal paavises (19). --- Det er dog stedse kun
en Gradsforskel, der er mellem primære og sekundære Domme.
Det ses af, at de kunne formuleres ganske paa samme Maade.

% --- p. 23 ---
\opage{23}\chaphead{III}{ANALYSE AF DOMME}{SLUTNINGSDANNELSE}


\sechead{A. De logiske Grundsætninger.}{\textbf{A. De logiske Grundsætninger.}}


21. Allerede ved Overgangen fra Anskuelse til Begreb og
Dom gør sig, hvad enten man er sig det bevidst eller ikke,
en Grundsætning gældende, nemlig den, at Begrebets eller
Dommens Indhold skal svare til Anskuelsens Indhold, saa at
hin Overgang kan betragtes som Omsætning til en anden
Form. Ved Begrebs- og Domsdannelse frembringes der i en
vis Forstand intet Nyt, uden netop selve den nye Form, der
bliver af Betydning for det fortsatte aandelige Arbejde. Men
endnu tydeligere fremtræder den samme Grundsætning ved
Slutningen, d. v. s. ved Dannelsen af nye Domme ud af givne
Domme. Slutsætningen (Konklusionen) skal svare til Forudsætningerne
(Præmisserne) og kun være en Omsætning af
deres Indhold i en ny Form. Forholdet mellem Slutsætning og
Forudsætninger er analogt det mellem Begreb eller Dom og
Anskuelse. Det er det Forhold, som kaldes Forholdet mellem
\textit{Grund og Følge}. Ved Slutningsdannelsen træder dette Forhold
tydeligt frem; derfor drøftes de logiske Grundsætninger som
Indledning til Slutningslæren, skønt de allerede kunne spores
i Læren om Begreb og Dom.

22. Den mest almindelige Grundsætning for Tænkningen
er, at vore Begreber og Domme skulle \textit{begrundes}, d. v. s. føres
tilbage enten til givne Anskuelser eller til andre Begreber og
Domme, som antages gyldige. Vore Begreber og Domme danne
vi fra først af uvilkaarligt, og Opgaven at begrunde dem kommer
bagefter. Logikens første Fordring er, at Begreber og
% --- p. 24 ---
% running head: III . Analyse af Domme .
\opage{24}Domme skulle have deres tilstrækkelige Grund: \textit{den tilstrækkelige
Grunds Princip} (principium rationis sufficientis). Hvor
langt vi skulle og kunne gaa tilbage ved Begrundelsen, beror
paa det enkelte Spørgsmaals Natur. Ofte maa vi stanse paa
Grund af manglende Data; men i sig selv fører det nævnte
Princip ind i en uendelig Række af Opgaver.

Om der i Virkeligheden er en Grund til Alt, hvad der sker,
er et Spørgsmaal, den formelle Logik, der kun har at gøre
med Ordningen af vore Tanker, ikke kan afgøre. Der maa
skelnes mellem Grund og Aarsag. Begreber og Domme have
deres Grund i Anskuelser; Slutsætninger have deres Grund
i Domme, der antages gyldige. Men Aarsagsforholdet antages
at finde Sted mellem Begivenheder (i os eller udenfor os),
der afløse hinanden i Tiden, og som ere saaledes knyttede til
hinanden, at naar den ene indtræder, indtræder uundgaaeligt
ogsaa den anden. Mellem Forudsætninger og Slutsætning er
der intet Tidsforhold; denne er allerede indeholdt i („ligger i“)
hine, selvom vi maaske ere længe om at indse det. Men ved
Aarsagsforholdet gaar den Begivenhed, der er Aarsag, i Tiden
forud for den, der er Virkning. Mellem Begreberne Grund og
Følge paa den ene Side, Begreberne Aarsag og Virkning paa
den anden Side er der en Analogi, som udtrykkes ved, at Aarsagen
kaldes Virkningens Tilblivelsesgrund (foruden at den for
os er Erkendelsesgrund, d. v. s. Grund til at vente den Begivenhed,
vi kalde Virkningen). Aarsagsproblemet, som det er
Erkendelsesteoriens, ikke den formelle Logiks Opgave at drøfte,
angaar netop Berettigelsen af Analogien mellem vore Sætningers
Begrundelse og Begivenhedernes Tidsfølge. --- I formel
Logik og ren Matematik have vi kun med Grunde at gøre;
i de reale Videnskaber (Naturvidenskab, Psykologi og Historie)
have vi tillige med Aarsager at gøre. Forskellen mellem de
formelle og de reale Videnskaber beror paa den Rolle, Kvalitetsforskelle
og Tidsforskelle spille i disse sidste, medens de formelle
% PRINTED AS IS: »Videnkaber« (for Videnskaber)
Videnkaber se bort fra dem. (For den formelle Logik
ere dog ogsaa Figurer og Størrelser Kvaliteter, smlgn. 2.)

23. Vor Stræben efter Begrundelse kan stanse paa to
Punkter, enten ved de Anskuelser, af hvilke vi danne de
% --- p. 25 ---
% running head: A. De logiske Grundsætninger .
\opage{25}simpleste Begreber og Domme, eller ved de Grundsætninger,
der forudsættes ved Dannelsen af Begreber og Domme og
ved Uddragen af Slutninger. I første Tilfælde maa det undersøges,
om vore Sanseanskuelser ikke ere illusoriske, om vore
Erindringer ikke bedrage os, og om vi kunne og ville fastholde
vore Fantasibilleder. I andet Tilfælde vil en fortsat Diskussion
være umulig, naar Nogen blandt Deltagerne ikke anerkender
de logiske Grundsætninger; disse kunne, da de forudsættes
ved al Begrundelse, ikke selv begrundes. Den, der vilde negte
al Begrundelses Forudsætninger, vilde stille sig selv udenfor
al Diskussion. Vi opdage de logiske Grundsætninger ved at
lægge Mærke til den Fremgangsmaade, der anvendes ved Dommes
Udledning af andre Domme eller af Anskuelser. Den, der
negter dem, maa enten paavise en anden Fremgangsmaade
eller opgive al Diskussion.

Ved Begrundelse af Domme og Opstilling af Begreber
kunne fem Principer være nødvendige.

24. \textsuperscript{1}) Identitetsprincipet: $A = A$. Dette Princip ligger allerede
i Definitionen af et Begreb som en fast og almengyldig
Forestilling, d. v. s. en Forestilling, der overalt, hvor den optræder,
har det samme Indhold, altsaa kan sættes lig sig selv.
Det ligger fremdeles i Forholdet mellem Subjekt og Prædikat
i Dommen, naar denne formuleres som Identitet. Og det forudsættes
ved al Slutten: af to Domme ($A = B$ og $B = C$) kan
jeg kun udlede en tredie ($A = C$), naar de have et fælles Begreb;
jeg kan da sætte Udtrykket for dette Begreb fra den
ene Dom ind i den anden. Slutningen sker ved Substitution,
og Identitet kan defineres som Substitutionsmulighed. Identiske
Begreber maa overalt kunne sættes i Stedet for hinanden. Jeg
har kun Ret at sætte $A = B$ (eller $A = AB$ o. s. v.), naar de
Begreber, der forbindes ved Lighedstegnet, kun ere forskellige
Udtryk for et og samme Indhold, saa at jeg overalt kan sætte
det ene i Stedet for det andet.

Vore virkelige Forestillinger kunne ikke nogensinde godtgøres
at være absolut identiske. Vi kunne kun paavise en faktisk
Dækningslighed, hvor der ingen Forskel mærkes (smlgn.
Psykologi VB, 8. 11); men absolut Sikkerhed for Identiteten
% --- p. 26 ---
% running head: III . Analyse af Domme .
\opage{26}kan ikke naas. Selv hvor vi tænke en og samme Tanke igen
(f. Ex. gaa fra $A$ paa den ene Side af Lighedstegnet til $A$ paa
den anden Side af Lighedstegnet), kan der være en Nuanceforskel
i den Maade, hvorpaa vi tænke det. Ved at opstille
Identitetsprincipet giver Logiken os en Maalestok eller et Ideal,
hvorved vi kunne prøve vore Begrebers, Dommes og Slutningers
Nøjagtighed og Gyldighed. Kun en Tilnærmelse i Retning
af dette Ideal kan være mulig.

25. \textsuperscript{2}) \textit{Enkelthedens Princip}: $AAA ..... = A$. Et Begrebs
Gentagelse og Kombination med sig selv har ingen logisk Betydning,
da dets Indhold ikke ændres derved. --- Medens Logiken
ikke kan faa Noget ud af identiske Begrebers Kombination,
er det anderledes i Aritmetiken: her lægges Vægt paa,
at de gentagne Begreber forudsætte ligesaa mange forskellige
Tankeakter, og at Talbegrebet derfor ligesaa vel finder Anvendelse
ved Rækken $AAA$ .... som ved Rækken $ABC$ .... Ved Talbegrebets
Anvendelse kan det være ligegyldigt, om Enerne have
samme eller forskelligt Indhold. --- For Psykologien er et Bevidsthedselements
Gentagelse af stor Betydning; den gør Genkendelse
mulig; den har Tendens til at svække Følelsen, men
kan paa den anden Side selv være Udtryk for stærk Følelse,
idet vi f. Ex. gentage Ordene, forat Følelsen ret kan faa Luft
(Sandeligt, sandeligt jeg siger Eder!). (Smlgn. Psykologi VI E,
4; F, 4.)

26. \textsuperscript{3}) \textit{Kontradiktionsprincipet} (Modsigelsens Grundsætning):
$Aa = 0$. Jeg kan ikke danne en Kombination eller en Dom
ved at forbinde et Begreb med dets egen Negation. Hvor en
saadan Forbindelse (f. Ex. „en rund Firkant“ eller „Firkanten
er rund“) er dannet, maa den strax ophæves. For Logiken er
en selvmodsigende Forbindelse et Intet (= 0). --- For Psykologien
ere de faktiske Selvmodsigelser Fænomener, der skulle
forklares. Selvmodsigelse er mulig, saalænge de modsigende
Forestillinger ikke opfattes bestemt, ikke analyseres nøjagtigt
og ikke bringes i umiddelbar Forbindelse med hinanden. Mangehaande
psykologiske Aarsager (især Følelsens Indflydelse, se
Psykologi VI F, 1) kunne hindre Bevidstheden i at opdage
Selvmodsigelsen.

% --- p. 27 ---
% running head: A. De logiske Grundsætninger .
\opage{27}27. \textsuperscript{4}) Medens jeg ikke kan kombinere $A$ med $a$, kan jeg
kombinere det med et andet Begreb. Og jeg har da, hvad
ethvert saadant Begreb angaar, Valget mellem to Muligheder:
at forbinde $A$ med dette Begreb (f. Ex. $B$) eller ikke forbinde
$A$ med det. Der er kun disse to Muligheder; vi staa altsaa
her overfor et kontradiktorisk Forhold. Der gives kun to logiske
Dimensioner: Position og Negation. (Smlgn. 18.) Dette
Forhold er udtrykt i \textit{Dualitetsprincipet} (principium exclusi
medii inter duo contradictoria): Begrebet $A$ maa enten forbindes
med $B$ eller ikke forbindes med $B$. Dette kan udtrykkes saaledes:
$AB \udel Ab$ (idet Tegnet $\udel$ betyder et gensidigt Udelukkelsesforhold,
et kontradiktorisk Forhold, et Enten---Eller).
Dualitetsprincipet udtrykker baade Uforeneligheden af de to
Kombinationer og det Forhold, at den enes Ophævelse er
Hævden af den anden.

Et Hvirveldyr maa enten være et Pattedyr eller ikke være
det: de Kendemærker, der ere ejendommelige for Hvirveldyrstypen,
maa enten forekomme sammen med dem, der ere ejendommelige
for Pattedyrstypen, eller uden disse.

Ethvert Begreb og dets Negation dele Verden i to Dele
($B$ og $b$). Hvad der ikke har det ene Begrebs Kendemærker,
maa have det andets. Og det Samme gentager sig indenfor
hvert enkelt Begreb, naar en Inddeling skal finde Sted. De
negative Begreber spille derfor en vigtig Rolle ved Forsøg
paa Klassifikation. (Smlgn. 8.)

Naar $B$ ikke er en nødvendig Bestemmelse ved $A$, kunne
baade $AB$ og $Ab$ forekomme --- under forskellige Forhold eller
til forskellige Tider. Hvirveldyrstypen forekommer baade hos
Pattedyr og hos Ikke-Pattedyr, f. Ex. baade hos Hesten og
hos Ørnen.

28. \textsuperscript{5}) \textit{Ombyttelighedens Princip}: $AB = BA$. (Smlgn. 14 b.)
Det er ligegyldigt, om jeg begynder med det ene eller med
det andet Begreb i en Kombination, ti Indholdet ændres ikke
derved. Leddenes Orden er ligegyldig. --- Samme Regel gælder
i Aritmetiken for Additions og Multiplikations Vedkommende.
Derimod er det i psykologisk Henseende (som overalt hvor
der er Tale om reale Begivenheder) ikke ligegyldigt, i hvilken
% --- p. 28 ---
% running head: III . Analyse af Domme .
\opage{28}Orden Leddene følge efter hverandre, da de foregaaende Led
faa Indflydelse paa de følgende. (Psykologi VA, 4; VI E.)

29. De logiske Principer ere ikke blot Forudsætninger for
al Begrundelse, og derved for al Diskussion, men de udtrykke
tillige den menneskelige Tænknings Natur. At tænke er at forbinde
gennem Sammenligning (se 1 og Psykologi VB, 11), og
Tænkningen bevæger sig frem gennem Identifikationer, Substitutioner,
Negationer og Kombinationer, hvis Grundregler ere
formulerede i de logiske Principer. Endeligt føre de logiske
Principer til Opstilling af Erkendelsesidealer og virke som
ledende Tanker ved vore Undersøgelser, idet de føre til at
opfatte baade Ligheder og Forskelle saa nøjagtigt som muligt
og derved gøre vore Tankeforbindelser stedse sikrere og fuldkomnere.
Om de gælde for Tilværelsen idethele, saa at denne,
set fra et menneskeligt Synspunkt, er rationel, er et Spørgsmaal,
der dels hører under Erkendelsesteorien, dels under Kosmologien
(Metafysiken). (Smlgn. Afsnittet „De filosofiske Problemer“
i „Kort Oversigt over den nyere Filosofis Historie“.)

\sechead{B. Slutninger.}{\textbf{B. Slutninger.}}


30. En Slutning gaar ud paa at bringe til tydelig Bevidsthed,
hvad der er indeholdt i visse givne Domme. Derved
frembringes intet nyt Indhold, men vi blive os bevidste, hvad
vi vide. Snart vide vi Mere, snart Mindre, end vi troede at
vide. Men det kan betegne et meget stort Fremskridt, at vi
blive os Betydningen og Rækkevidden af vore Domme bevidst.

Slutningen kan være af to Arter. Ved \textit{direkte Slutning}
haves to eller flere Domme givne, og det gælder da at finde
et fælles Begreb (Mellembegrebet), saaledes at der kan finde
Substitution Sted fra den ene Dom til den anden. Her skrides
altsaa ved Hjælp af Identitetsprincipet lige frem fra Forudsætninger
til Slutsætning. Ved \textit{indirekte Slutning} undersøges,
hvilke Kombinationer af de i en eller flere Domme indeholdte
Begreber, der kunne bestaa med disse Dommes Gyldighed.
Denne Undersøgelse sker ved Forkastelse af de Kombinationer,
der ved Substitution ud fra de givne Domme vise sig selv%
% --- p. 29 ---
% running head: B. Slutninger .
\opage{29}modsigende. Her bygges altsaa ikke blot paa Identitetsprincipet,
men ogsaa paa Kontradiktionsprincipet (og som det siden vil
vise sig, Dualitetsprincipet), og Resultatet naas gennem Udskydning
af de ugyldige Kombinationer. --- Ved begge Arter
af Slutning kunne Enkelthedens og Ombyttelighedens Principer
finde Anvendelse.

Slutningen beror, ligesom Begrebsdannelsen og Domsdannelsen,
paa Vexelvirkning af Analyse og Syntese. De givne Domme,
der ere opstillede og sammenstillede ved en Proces, der svarer
til den første Syntese (1. 4. 10), maa analyseres, forat man
kan finde, hvilke Begreber der kunne substitueres. Men dernæst
finder ved Hjælp af Substitution en ny Syntese Sted,
som fører til at se, hvilke Domme eller Kombinationer der
svare til og bestaa med de givne Domme. Slutningen er, som
allerede Aristoteles har udtrykt det, en \textit{Sammentænken} af det
Givne. Det græske Ord for Slutning (syllogismos) betyder
netop Sammentænken. Slutningen er en logisk Syntese.

\sechead{a. Direkte Slutning.}{\textbf{a. Direkte Slutning.} (Slutning ved Substitution alene.)}


31. Der forudsættes her, at vi have (i simpleste Tilfælde)
to Domme givne, der have et fælles Begreb, saa at Substitution
fra den ene Dom i den anden bliver mulig. Et saadant Begreb
kaldes Mellembegrebet; de Begreber, i Forhold til hvilke Mellembegrebet
enten er Subjekt eller Prædikat, kaldes Hovedbegreberne.
--- Af de utallige Former for direkte Slutning
anføres her kun nogle enkelte som Exempler.

32. \textit{Naar et Subjekt har et Prædikat, der igen er Subjekt
for et andet Prædikat, sluttes, at hint Subjekt ogsaa har
dette andet Prædikat}:
% scheme n2r-01-01 (notatex.py)
\[\begin{array}{l}
A = AB \\
B = BC \\
\hline
A = ABC = AC
\end{array}\]
Neptun er en Planet; enhver Planet bevæger sig i en Ellipse;
altsaa bevæger Neptun sig ogsaa i en Ellipse. Et Kvadrat
er et Parallelogram; i et Parallelogram halvere Diago%
% --- p. 30 ---
% running head: III . Analyse af Domme .
\opage{30}nalerne hinanden; altsaa halvere Diagonalerne i et Kvadrat
hinanden.

Denne Slutning udtrykkes i den aristoteliske Logik, der
bygger paa Begrebernes Omfang (se 13), saaledes:
% scheme n2r-01-02 (notatex.py)
\[\begin{array}{l}
A < B \\
B < C \\
\hline
A < C.
\end{array}\]
Den kaldes en Slutning i \textit{første Figur}. Aristoteles skelner
mellem forskellige „Figurer“ (Skemaer), idet efter hans Fremstilling
Begrebernes forskellige Plads som Subjekt eller Prædikat
i Præmisserne blev af stor Vigtighed. I første Figur er
Mellembegrebet Prædikat i den ene Præmis, Subjekt i den
anden. --- Ved geometrisk Fremstilling vil $A$ være en lille
Cirkel indeni $B$, og $B$ igen en Cirkel indeni $C$, og det ses
da umiddelbart, at $A$ indeholdes i $C$. ---

Hvis vi istedetfor partielle Identiteter havde absolute Identiteter
i Præmisserne, vilde Slutsætningen ogsaa være en absolut
Identitet:
% scheme n2r-01-03 (notatex.py)
\[\begin{array}{l}
A = B \\
B = C \\
\hline
A = C.
\end{array}\]
Denne Slutningsform, (der ikke kunde forekomme hos Aristoteles
og ikke kan fremstilles geometrisk) benyttes, naar Fænomeners
Identitet skal godtgøres. Har jeg en nøjagtig Beskrivelse
($B$) af en Mand ($A$), og denne Beskrivelse ganske
passer paa den Person ($C$), der begik et vist Mord, slutter
jeg, at $A$ er Morderen. Den Kraft ($A$), der holder Maanen i
dens Bane, er den samme Kraft ($C$), der paa Jorden kaldes
Tyngde, fordi Loven ($B$) for begge Kræfter er den samme. ---

Har jeg
% scheme n2r-01-04 (notatex.py)
\[\begin{array}{l}
A = AB \\
B = BA \\
\hline
\end{array}\]
følger $A = B$.

% --- p. 31 ---
% running head: B. Slutninger .
\opage{31}Her anvendes foruden Identitetsprincipet ogsaa Ombyttelighedens
Princip. En absolut Identitet sluttes her af to partielle
Identiteter. Det er egentligt kun en Definition af absolut
Identitet, vi her faa. (Smlgn. 14 b og 28.) --- Et ligevinklet Triangel
er ligesidet, og et ligesidet Triangel er ligevinklet; de to
Arter af Triangler ere altsaa identiske.

33. \textit{Deraf, at to Subjekter have et og samme Prædikat,
kan Intet sluttes om Forholdet mellem dem}.
% scheme n2r-01-05 (notatex.py)
\[\begin{array}{l}
A = AB \\
C = CB
\end{array}\]
Her er Intet at substituere. Hos Aristoteles bliver dette
\textit{anden Figur}:
% scheme n2r-01-06 (notatex.py)
\[\begin{array}{l}
A < B \\
C < B,
\end{array}\]
og han advarer udtrykkeligt mod at slutte her, med mindre
den ene Præmis er negtende. (Naar $A$ er $B$, men $C$ ikke er
$B$, kan $A$ ikke være $C$.) --- Geometrisk ligge baade $A$ og $C$
indenfor Cirklen $B$; men deres indbyrdes Forhold kan ikke
afgøres; de kunne dække hinanden eller skære hinanden eller
falde udenfor hinanden eller den ene indenfor den anden.
--- Deraf, at baade Fisk og Hvaler leve i Vandet, kan hverken
sluttes, at Hvaler ere Fisk, eller omvendt. Men deraf, at Fisk
have koldt Blod og Hvaler varmt Blod, kan sluttes, at Hvaler
ikke ere Fisk. Hvis $A = AB$ og $C = Cb$, kan $C$ ikke være
Prædikat til $A$, da vi saa vilde faa $ABCb$, der indeholder en
Modsigelse.

34. \textit{Naar et og samme Subjekt har to Prædikater, kan
deraf sluttes en begrænset Identitet af de to Prædikatsbegreber}.
% scheme n2r-02-01 (notatex.py)
\[\begin{array}{l}
B = BA \\
B = BC \\
\hline
BA = BC
\end{array}\]
Fiskene ere baade Vanddyr og Hvirveldyr, og disse to
Egenskaber forekomme da under visse Betingelser sammen.
% --- p. 32 ---
% running head: III . Analyse af Domme .
\opage{32}Mennesket er et aandeligt Væsen, og Mennesket er et materielt
Væsen; det Aandelige og det Materielle maa da være
foreneligt under visse Betingelser. (Om de Betingelser, der
udtrykkes ved $B$, ere de eneste, under hvilke $A$ og $C$ kunne
forekomme sammen, siges her intet om.)

Denne Slutning fremstilles i den aristoteliske Logik saaledes:
% scheme n2r-02-02 (notatex.py)
\[\begin{array}{l}
B < A \\
B < C
\end{array}\]
og kaldes en Slutning i \textit{tredie Figur}. Mellembegrebet er her
Subjekt i begge Præmisser. Slutsætningen bliver: nogle $A$ er
$C$ (f. Ex.: nogle Vanddyr ere Hvirveldyr). --- Geometrisk kan
tredie Figur fremstilles ved, at Cirklerne $A$ og $C$ skære hinanden,
og at $B$ er det Stykke, der hører til begge Cirkler.
(Hvis $A < C$, komme vi tilbage til første Figur.) ---

Dersom vi ikke blot have $B = BA$ og $B = BC$, men ogsaa
$A = AB$ og $C = CB$ (altsaa kunne forbinde 2. og 3. Figur),
faa vi $A = C$, saa at de to Begreber, der først kun syntes at
være begrænset identiske, nu vise sig at være absolut identiske.
Hvis $A$ betyder „højst udviklet Fornuftvæsen“, $B$ „Menneske“
og $C$ „Dyr med fuldkomnest Hjerne“, faa vi, at det
højest udviklede Fornuftvæsen er det Dyr, der har den fuldkomneste
Hjerne.

35. \textit{Dersom begge Præmisser ere negtende}, kan man finde
en Slutsætning ved Benyttelsen af, hvad der i 18 er fremsat
om negtende Domme.

Har jeg de to nægtende Domme $A \sim B$ og $B \sim C$, maa jeg
undersøge, om de hver for sig blot ere kontrære, eller om de
ere kontradiktoriske. Kontrære maa enhver af dem nødvendigvis
være.

$A \sim B$ maa betyde \textsuperscript{1}) $B = Ba$, men kan ogsaa betyde \textsuperscript{2}) $B = a$.

$B \sim C$ maa betyde \textsuperscript{3}) $B = Bc$, men kan ogsaa betyde \textsuperscript{4}) $B = c$.

Hvis kun \textsuperscript{1}) og \textsuperscript{3}) gælde, faas $Ba = Bc$. Hvis \textsuperscript{1}) og \textsuperscript{4}) gælde,
faas $Ba = c$. Hvis \textsuperscript{2}) og \textsuperscript{3}) gælde, faas $a = Bc$. Hvis \textsuperscript{2}) og \textsuperscript{4})
gælde, faas $a = c$.

Her er gaaet ud fra Mellembegrebet; men man kan ogsaa
gaa ud fra Hovedbegreberne:

% --- p. 33 ---
% running head: B. Slutninger .
\opage{33}$A \sim B$ maa betyde \textsuperscript{1}) $A = Ab$, men kan ogsaa betyde \textsuperscript{2}) $A = b$.

$B \sim C$ maa betyde \textsuperscript{3}) $C = Cb$, men kan ogsaa betyde \textsuperscript{4}) $C = b$.

Hvis \textsuperscript{1}) og \textsuperscript{3}) gælde, kan Intet sluttes (da det er anden Figur).
Hvis \textsuperscript{1}) og \textsuperscript{4}) gælde, faas $A = AC$. Hvis \textsuperscript{2}) og \textsuperscript{3}) gælde,
faas $C = CA$. Hvis \textsuperscript{2}) og \textsuperscript{4}) gælde, faas $A = C$.

36. Undertiden er der kun én Forudsætning udtrykkeligt
given, og dog kan en Slutning drages. En saadan Slutning
kaldes en \textit{umiddelbar Slutning}. Den staar som selvindlysende.
Men fordres der et Bevis, behøves dertil et eller flere af de
logiske Principer. Her (ved direkte Slutning) have vi kun at
gøre med de umiddelbare Slutninger, der kunne bevises ved
Identitetsprincipet.

Dersom $A = B$, følger $A = AB$ og $B = BA$. Dette ligger
og allerede i selve Begrebet absolut Identitet. (Smlgn. 14 b. 28. 32.)
Men det kan bevises saaledes: $A = AA$ (efter Enkelthedens
Princip). Substitueres nu $B$ et Sted paa højre Side, faas $A = AB$.
Paa samme Maade faas af $B = BB$, at $B = BA$. --- Af $A = B$
kan ogsaa sluttes $AC = BC$. Hvad der gælder absolut, maa
ogsaa gælde under en vis Betingelse. Beviset er følgende: Vi
sætte: $AC = AC$ og indsætte paa den ene Side $B$ for $A$, saa
faas $AC = BC$.

37. \textit{Det Spørgsmaal, om en Slutning kan vendes om}, saa
at Slutsætningen bliver Præmis, og de tidligere Præmisser Slutsætninger,
maa undersøges i de forskellige Tilfælde.

% PRINTED AS IS: the cross-reference »(33)«; § 32 may be meant (the inference from A = AB and B = BA to A = B is set out there)
Slutningen fra $A = AB$ og $B = BA$ til $A = B$ (33) kan
vendes om (se 36). Men Slutningen fra $A = B$ til $AC = BC$
(36) kan ikke vendes om, da man ikke uden videre har Ret
til at ophæve Betingelsen. --- Hvor Mellembegrebet forsvinder
i Slutningen, som naar der (32) sluttes fra $A = AB$ og $B = BC$
til $A = AC$, eller fra $A = B$ og $B = C$ til $A = C$, kan Slutningen
ikke vendes om, da man ikke af Slutsætningen kan
finde Mellembegrebet, og da desuden ethvertsomhelst Mellembegreb
vil give samme Slutsætning. $A = C$ følger ligesaa vel
af $A = X$ og $X = C$, som af $A = B$ og $B = C$.

Grunden (Forudsætningerne) fører med Nødvendighed til
Følgen (Slutsætningen), men kan ikke altid gaa den modsatte
Vej. I Erfaringsvidenskaben sluttes tilbage fra et givet Fæno%
% --- p. 34 ---
% running head: III . Analyse af Domme .
\opage{34}men (f. Ex. Tilstedeværelsen af Kulsyre) til Betingelserne for dets
Opstaaen; men disse kunne i forskellige Tilfælde være forskellige.
(Kulsyre kan opstaa ved Dyrenes Aandedræt, ved
Forbrænding, ved vulkanske Udbrud o. s. v.) Derfor ere de
fleste Naturforklaringer Hypoteser, begrundede Formodninger.
Og en Hypotese bliver ikke fuldt ud sikker, fordi den bekræftes
ved Erfaring, saalænge det ikke kan godtgøres, at
ingen andre Forudsætninger kunde føre til samme Resultat,
--- saalænge altsaa Slutningen ikke kan vendes om. Allerede
Leibniz har indset dette: le succès ne démontre pas la vérité
de l’hypothèse. Hele den nyere Naturvidenskab hviler saaledes
% PRINTED AS IS: the semicolon before »der opstilledes«
paa Forudsætninger; der opstilledes af Galilei og hans
Efterfølgere, og Spørgsmaalet er, om disse Forudsætninger ere
de eneste, ud fra hvilke de vundne Resultater kunne naas.
(Smlgn. herom de nyere, i „Oversigt over den nyere Filosofis
Historie“ VIII, 2 omtalte erkendelsesteoretiske Synsmaader.)

Ikke blot kan en og samme Slutsætning udledes af Domme
med højst forskellige Mellembegreber, men en rigtig Slutsætning
kan udledes af faktisk urigtige eller endog selvmodsigende
Præmisser. Baade $A = B$ og $B = C$ kunne være faktisk urigtige,
og dog kan $A = C$ være faktisk rigtig. Det viser sig ofte,
at man fra falske Forudsætninger kommer til et Resultat, der
saa ad andre Veje bevises at være rigtigt. Den rigtige Udledning
af $A = C$ kunde f. Ex. være den af $A = K$ og $C = K$. ---
Af $A = a$ og $a = C$ følger $A = C$, der ingen Selvmodsigelse
indeholder, skønt den ene Præmis er selvmodsigende.

\sechead{b. Indirekte Slutning.}{\textbf{b. Indirekte Slutning.} (Slutning ved Substitution og Elimination.)}

38. Den indirekte Slutning undersøger, hvilke Begrebskombinationer
der kunne bestaa med Gyldigheden af en eller flere
givne Domme. Denne Undersøgelse sker ved, at man i de
(ifølge Dualitetsprincipet) tænkelige Kombinationer af de i Dommen
(eller Dommene) givne Begreber indsætter de Udtryk,
Dommene give for disse Begreber og udskyder (eliminerer)
de Kombinationer, der blive selvmodsigende. Dette er en Deduktion,
idet man gaar fra en eller flere givne Domme til de
Kombinationer, der kunne bestaa med dem. --- De Domme,
% --- p. 35 ---
% running head: B. Slutninger .
\opage{35}som undersøges paa denne Maade, kunne selv være udledte
ved \textit{Induktion} af visse givne Begrebskombinationer. Deduktionen
giver da Beviset for, at Dommene ere rigtigt udledte.
Induktionen er væsentligt en Opdagelsesmetode; vi skaffe os
ved den Domme, der kunne prøves deduktivt. Den er et forberedende
Arbejde; den egentlige Slutning er altid deduktiv.

\sechead{α. Indirekte Slutning som ren Deduktion.}{\spaced{α. Indirekte Slutning som ren Deduktion.}}


39. Naar Dommen $A = B$ er given, hvilke Kombinationer
af $A$ og $B$ og deres Negationer ville da være mulige? --- Da
der er Tale om to Begreber, kan der \textit{ifølge Dualitetsprincipet}
kun tænkes fire Kombinationer: $A$ kan enten forekomme med
$B$ eller uden $B$; $B$ enten med $A$ eller uden $A$; $a$ enten med
$B$ eller uden $B$; $b$ enten med $A$ eller uden $A$. De fire Kombinationer,
der kunne tænkes, ere altsaa $AB$, $Ab$, $aB$ og $ab$.
(Geometrisk kan dette anskueliggøres ved at lade Cirklerne
$A$ og $B$ skære hinanden; der bliver da en Del af $A$ uden $B$,
en Del, som baade hører til $A$ og $B$ og en Del af $B$ uden $A$;
dertil kommer saa Rummet udenfor de to Cirkler, hvor hverken
$A$ eller $B$ findes.) Men nu gælder det at finde, hvilke af
disse tænkelige Kombinationer der ere mulige, naar Dommen
$A = B$ er given. Det sker ved at indsætte denne Dom i de
fire Kombinationer \textit{(i Kraft af Identitetsprincipet)} og derpaa
\textit{(i Kraft af Kontradiktionsprincipet)} udskyde de Kombinationer,
der blive selvmodsigende. De Kombinationer, i hvilke
Intet kan substitueres, blive staaende som mulige.
% scheme n2r-02-03 (notatex.py)
\[\begin{array}{l}
AB = AA = BB \\
Ab = Bb = 0 \\
aB = aA = 0 \\
ab
\end{array}\]
Altsaa kun $AB$ og $ab$ kunne forekomme, naar $A = B$. ---
Hvis Personen $A$ skal være identisk med Morderen $B$, saa
maa han have været tilstede overalt, hvor Morderen har været,
og have været borte, hvor Morderen var borte.

40. Naar Dommen $A = AB$ er given, faa vi efter samme
Fremgangsmaade $AB$, $aB$ og $ab$ som mulige Tilfælde. --- Hvis
% --- p. 36 ---
% running head: III . Analyse af Domme .
\opage{36}Varme er en Form for Bevægelse, maa vi kunne finde Varme- og
Bevægelsesfænomener sammen, eller Bevægelse uden Varme
(da der er andre Bevægelsesformer end Varme), eller Fænomener,
der hverken ere Bevægelse eller Varme. Kun Kombinationen
Varme uden Bevægelse ($Ab$) er udelukket.

Naar $A = Ab$ er given, blive $Ab$, $aB$ og $ab$ mulige. Naar
$A = b$ er given, blive kun $Ab$ og $aB$ mulige.

41. For Simpelheds Skyld er der i det Foregaaende valgt
Exempler med to Begreber. Hvis der er Tale om tre Begreber,
faa vi 8 tænkelige Kombinationer at undersøge. (Geometrisk
kan man fremstille dette ved at lade tre Cirkler skære
hverandre saaledes, at enhver af dem skærer begge de andre;
man faar da 7 forskellige Stykker foruden Rummet udenfor
de tre Cirkler.) Er der fire Begreber indeholdte i de givne
Domme, faa vi 16 Kombinationer o. s. v. --- Følgende Exempel
med tre Begreber viser Fremgangsmaaden.

Lad $A = AB$ og $C = CB$ være givne. Saa ere de 8 tænkelige
Kombinationer:
% scheme n2r-02-04 (notatex.py)
\begin{center}
\begin{tabular}{l@{\qquad}l}
\textsuperscript{1}) $ABC$ & \textsuperscript{5}) $aBC$ \\
\textsuperscript{2}) $ABc$ & \textsuperscript{6}) $aBc$ \\
\textsuperscript{3}) $AbC$ & \textsuperscript{7}) $abC$ \\
\textsuperscript{4}) $Abc$ & \textsuperscript{8}) $abc$
\end{tabular}
\par\smallskip Ved Substitution faas:\par\smallskip
\begin{tabular}{l@{\qquad}l}
\textsuperscript{1}) $ABBCB = ABC$ & \textsuperscript{5}) $aBCB = aBC$ \\
\textsuperscript{2}) $ABBc = ABc$ & \textsuperscript{6}) $aBc$ \\
\textsuperscript{3}) $ABbCB = 0$ & \textsuperscript{7}) $abCB = 0$ \\
\textsuperscript{4}) $ABbc = 0$ & \textsuperscript{8}) $abc$
\end{tabular}
\end{center}
Man kan altsaa ved indirekte Slutning drage Slutninger fra
„anden Figur“ (33) angaaende de Begrebskombinationer, der
ere mulige, naar to Subjekter have samme Prædikat.

Gaar man ud fra $B = BA$ og $B = BC$ (altsaa „tredie Figur“,
se 34), blive $ABC$, $AbC$, $Abc$, $abC$ og $abc$ mulige.

42. Dersom vi gaa ud fra indbyrdes modsigende Præmisser,
vil indirekte Slutning kunne gøre opmærksom paa Fejlen.

% --- p. 37 ---
% running head: B. Slutninger .
\opage{37}%
% scheme n2r-02-05 (notatex.py)
\[\begin{array}{l@{\qquad}l}
A = b \\
B = A \\
\hline
AB = bB = 0 & aB = aA = 0 \\
Ab = bB = 0 & ab = aA = 0
\end{array}\]
Alle tænkelige Kombinationer falde altsaa bort. --- Et andet
Exempel vil forekomme nedenfor (46).

\whiteline

\sechead{β. Indirekte Slutning som Forbindelse af Induktion og Deduktion.}{\spaced{β. Indirekte Slutning som Forbindelse af Induktion}\\
\spaced{og Deduktion.}}

43. Have vi en Række Begrebskombinationer givne, kunne
vi søge en eller flere Domme, der udtrykke Forholdet mellem
de i Kombinationerne indeholdte Begreber. Saadanne Domme
findes ved Analyse, ligesom naar vi danne vore Domme paa
Grundlag af umiddelbar Anskuelse (10). Naar Kombinationerne
ere meget talrige, maa man prøve sig frem gennem foreløbige
Hypoteser. Saaledes prøvede Kepler først den Antagelse, at
Marsbanen var en Cirkel, før han opdagede, at den var en
Ellipse. Den hjælpende Syntese spiller her en stor Rolle.
Kombinationerne staa her jævnsides, medens det ved de simplere
Tilfælde af Domsdannelse gælder om at belyse en enkelt
(i Anskuelsen given) Kombination ved Hjælp af de andre.
(Smlgn. 10 og 20.)

Overgangen fra givne Kombinationer til Domme kaldes
Induktion. Men et Bevis for, at de fundne Domme virkeligt
svare til de givne Kombinationer og ikke til andre, maa føres
ved Deduktion. Beviset for en Induktion er stedse en Deduktion,
og Induktionens Betydning bestaar i at opstille en Formodning,
der kan prøves. (Smlgn. 38.)

Indirekte Slutninger maa stedse kunne vendes om, hvis de
ere rigtige. Jeg maa baade kunne gaa fra en bestemt Dom til
bestemte Kombinationer, og omvendt.

44. Ere Kombinationerne $AB$ og $ab$ givne, se vi, at hvor
$A$ er, er $B$, og omvendt, og at hvor en af dem mangler,
mangler ogsaa den anden. Dertil svarer $A = B$ og $a = b$. Og
ved Deduktion fra disse Domme komme vi ogsaa til $AB$ og
$ab$, og til dem alene (se 39).

% --- p. 38 ---
% running head: III . Analyse af Domme .
\opage{38}Ere $AB$, $aB$ og $ab$ givne, se vi, at hvor $A$ er, er $B$, medens
$B$ kan forekomme uden $A$, og at hvor $B$ mangler, mangler $A$.
Det giver $A = AB$ og $b = ba$. Om $B$ og $a$ kan Intet udsiges.
Deduktion fra $A = AB$ fører ogsaa til $AB$, $aB$ og $ab$, medens
$Ab$ vilde være selvmodsigende. (Se 40.)

Ere $AB$, $Ab$ og $aB$ givne, ses, at $a$ kun forekommer, hvis
$B$ er, medens $B$ kan forekomme uden $a$, og at hvor $B$ mangler,
mangler ogsaa $a$. Dertil svarer $a = aB$ og $b = bA$. Om $A$ og
$B$ kan her Intet udsiges. Ved Deduktion fra hine to Domme
faa vi ogsaa de tre givne Kombinationer og ikke andre.

Hvis de givne Kombinationer var $ABC$, $AbC$, $aBc$ og
$abC$, vilde deri ligge $A = AC$, $b = bC$, $c = cB$, fremdeles
$AB = ABC$ o. s. v.

45. Mange \textit{umiddelbare Slutninger} (smlgn. 36) kunne kun
bevises ved indirekte Slutning. Det gælder om den vigtige Art
af umiddelbar Slutning, ved hvilken vi \textit{af en positiv Dom
udlede en negativ Dom eller omvendt}. En saadan Afledning
kaldes \textit{Kontraposition}.

Af $A = AB$ følger umiddelbart $b = ba$. Ti naar $B$ er, hvor
$A$ er, maa $A$ mangle, hvor $B$ mangler. Bevis kan dog føres
saaledes: Af de fire mulige Kombinationer falde, naar $A = AB$
er given, $Ab$ bort; kun $AB$, $aB$ og $ab$ foreligge altsaa, og af
dem følger (se 44) $b = ba$.

Af $A = B$ følger umiddelbart $a = b$. Er der absolut Identitet
mellem $A$ og $B$, maa $a$ mangle, hvor $b$ mangler, og omvendt.
Bevis kan føres saaledes: Af de fire mulige Kombinationer
falde, naar $A = B$ er given, $Ab$ og $aB$ bort; kun $AB$ og $ab$
foreligge, og af dem følger (se 44) $a = b$.

Disse Slutninger kunne vendes om. Af $b = ba$ følger ved
Kontraposition $A = AB$, af $a = b$ følger $A = B$. ---

Formedelst Kontraposition kunne vi formulere en Slutning
i „anden Figur“ med negtende Præmis (se 33):

% CHECK p. 38: display of notation
\begin{center}$A = AB$\end{center}


% CHECK p. 38: display of notation
\begin{center}$C = Cb$\end{center}


Af $C = Cb$ følger $B = Bc$, saa at vi forvandle anden Figur
til første Figur:

% --- p. 39 ---
% running head: B. Slutninger .
\opage{39}%
% scheme n2r-03-01 (notatex.py)
\[\begin{array}{l}
A = AB \\
B = Bc \\
\hline
A = ABc = Ac.
\end{array}\]
Selv af anden Figur med to positive Præmisser kunne vi
nu faa et Resultat. Ti af $A = AB$ følger $b = ba$, og af $C = CB$
følger $b = bc$, saa at vi faa:
% scheme n2r-03-02 (notatex.py)
\[\begin{array}{l}
b = ba \\
b = bc \\
\hline
ba = bc
\end{array}\]
Anden Figur forandres altsaa ved Kontraposition til tredie
Figur. ---

46. Ved enhver Induktion maa vi have alle i Kombinationerne
indeholdte Begreber givne baade i negativ og i positiv
Form, da ethvert Begreb maa kunne negeres. (Smlgn. 5---6.)
Skal der sluttes fra to Kombinationer alene, maa disse altsaa
være enten $AB$ og $ab$ eller $Ab$ og $aB$ (naar der kun er Tale
om to Begreber).

Hvis vi derimod f. Ex. havde $AB$ og $aB$ givne, vilde vi
faa de to Domme $A = AB$ og $a = aB$. Men af dem følger ved
Kontraposition (45) $b = ba$ og $b = bA$, der stride mod hinanden
(gøre hinanden problematiske, se 19). Der opstaar altsaa
en Modsigelse, naar vi slutte fra givne Kombinationer, der
udelukke et Begrebs Negation. Denne Modsigelse kan kun
skaffes afvejen, naar vi kunne finde flere Kombinationer. Finde
vi f. Ex. $ab$, falder Dommen $b = bA$ bort, og Modsigelsen ophører.
Saaledes føre de Problemer, Erfaringen stiller, til at udvide
vor Erfaring. --- Omvendt vilde vi ved deduktivt at gaa
ud fra de to Sætninger $A = AB$ og $a = aB$ komme til det Resultat,
at Begrebet $B$ ikke skulde kunne negeres, og dette Resultat
vilde naturligt føre os til at opdage Modsigelsen mellem
Præmisserne. ---

Dersom alle tænkelige Kombinationer ere givne, kan ingen
Dom opstilles. Alle Begreber forekomme da i alle de i Kraft
af Dualitetsprincipet tænkelige Forbindelser: $A$ snart med $B$,
snart med $b$, $B$ snart med $A$, snart med $a$ o. s. v. Hvis vor
% --- p. 40 ---
% running head: III . Analyse af Domme .
\opage{40}Erfaring stedse var saaledes beskaffen, vilde vi ingen Domme
kunne danne, og hvis hele Naturen var saaledes beskaffen,
vilde Aarsagssætningen ikke gælde, da den forudsætter, at et
Fænomen kun indtræder i Forbindelse med et vist bestemt
andet Fænomen. Hvor der i Erfaringen synes at foreligge Tilfælde,
der frembyde alle tænkelige Kombinationer (ved to Begreber
altsaa baade $AB$, $Ab$, $aB$ og $ab$), antage vi --- hvis det
er os om at gøre at fastholde Aarsagssætningen --- at vor Erfaring
er ufuldstændig, idet den er begrænset til to Begreber.
Hvis vor Erfaring udvides saaledes, at vi tage Hensyn til et
tredie Begreb $C$, der forekommer sammen med $A$ og $B$, ville
vi maaske kunne opstille en eller flere Domme. Saaledes hvis
vi fandt Kombinationerne $ABC$, $AbC$, $aBc$ og $abC$ (se 44 Slutning).
Dersom derimod alle 8 Kombinationer af $A$, $B$ og $C$
forekom, kunde heller ikke her nogen Dom opstilles. Vi
maatte da opstille et fjerde Begreb $D$, der kunde høre sammen
med de tre Begreber. Og saaledes fremdeles, til Vanskeligheden
er hævet. (Smlgn. 14. En begrænset Identitet kan ikke
udledes af Subjekts- og Prædikatsbegrebet alene, men findes
ved en Udvidelse, der godtgør, at Forholdet mellem de to Begreber
kun kan fastsættes, naar et tredie Begreb tages med.)
--- Saaledes fører den tænkende Bearbejdelse af det Givne til
at se dets Ufuldstændighed og driver derved til at udvide Erfaringen.


47. Selvom den fundne Dom ved Deduktion fører til de
givne Kombinationer og ikke til andre, og altsaa viser sig at
være rigtigt udledet, vedbliver den dog at være en \textit{Hypotese},
saalænge det ikke med Sikkerhed er godtgjort, at de givne
Kombinationer ere udtømmende, saa at ingen Kombinationer
af anden Art forekomme. En saadan Sikkerhed kan, afset fra
meget simple og overskuelige Forhold, være meget vanskelig
at naa. Hypotesens videnskabelige Karakter beror paa, at den
er det logiske Udtryk for alle givne Kombinationer. Da staar
den som en \textit{begrundet} Formodning. Men den trænger til \textit{Bekræftelse}
(Verifikation), saalænge der endnu er Mulighed for
nye Kombinationer. Bekræftelsen sker da ved, at de nye Kombinationer
paavises at være af samme Art som de, til hvilke
% --- p. 41 ---
% running head: B. Slutninger .
\opage{41}man støttede sig ved Opstillingen af Hypotesen. (Smlgn. 5 om
Definitionen som en Hypotese.) Hypotesen $A = B$ bekræftes,
dersom alle nye Erfaringer om $A$ og $B$ ere af Arterne $AB$
eller $ab$; den gendrives, dersom Tilstedeværelsen af $Ab$ eller
$aB$ kan paavises. (Se desuden 37.)

Den videnskabelige Induktion gaar ud fra visse Tilfælde
som typiske og paaviser, at alle andre Tilfælde høre til samme
Type. Newtons Gravitationsteori begrundedes ved en Undersøgelse
af Maanens Bevægelse; den bekræftedes ved, at Keplers
Love kunde udledes af den, --- at Ebbe og Flod kunde
forklares ved den, --- og senere ved en Mængde nye Opdagelser,
f. Ex. Opdagelsen af Neptun, der fremkaldtes ved
Slutning fra Gravitationsteorien. En Hypoteses Arbejdsværdi
(dens Betydning som Princip for Forskning) falder sammen
med Muligheden af dens Verifikation.

% --- p. 42 ---
\opage{42}\chaphead{}{TILLÆG}{}


\sechead{1. Mill’s induktive Metoder.}{\textbf{1. Mill’s induktive Metoder.}}


JEVONS har først med Klarhed paavist, at Induktionsbeviset
bestaar i en Deduktion, der ud fra foreløbigt opstillede Domme
fører tilbage til de Kombinationer, af hvilke disse Domme udledtes,
i Forbindelse med en fortsat Godtgørelse af, at alle
Kombinationer ere af samme Art som de, af hvilke Dommene
udledtes. ARISTOTELES forstod ved Induktion en fuldstændig
Samling af en vis Art Fænomener med Paavisning af Domme,
der gælde for dem alle. Men en saadan fuldstændig Samling
er meget sjeldent mulig, og Problemet er netop, hvorledes vor
Erkendelse kan bygges op og føres videre, skønt det Givne
ikke kan bevises at være udtømmende. BACON og STUART
MILL have givet specielle Beskrivelser af Induktionsmetoderne,
uden at de have set deres nøje Sammenhæng med Deduktionen,
der ene kan give et Bevis. Stuart Mill’s Fremstilling,
der er bygget paa et Studium af den nyere Erfaringsvidenskabs
Historie, har dog stadigt Interesse ved den Klarhed, med hvilken
den fremhæver de væsentlige Former.

a. \textit{Overensstemmelsesmetoden}. Naar et Fænomen $B$ indtræder
under forskellige Omstændigheder, saaledes at der kun
er én Ting ($A$), der stedse gaar forud, medens alt Andet skifter,
formodes, at $A$ er Aarsag til $B$ (eller at begge ere succesive Virkninger
af samme Aarsag). Denne Formodning vil være des
mere begrundet, jo flere Tilfælde vi have af denne Art.
% scheme n2r-03-03 (notatex.py)
\[\begin{array}{l@{\qquad}l@{\qquad}l}
AVX & AYZ & ......... \\
BV_1X_1 & BY_1Z_1
\end{array}\]
% --- p. 43 ---
% running head: 1. Mill's induktive Metoder .
\opage{43}Sammenlignes de forskellige Tilfælde, i hvilke Dugdannelsen
indtræder (ved Overflader i fri Luft efter Solens Nedgang, paa
et Glas med frisk Vand i en varm Stue, paa Indersiden af
et Vindu ved pludseligt Regnskyl o. s. v.), frembyde de trods
alle Forskelligheder det fælles Træk, at der forud for Dugdannelsen
indtræder en pludselig Afkøling af en Overflade i
Forhold til den omgivende Luft, og deri finde vi da Aarsagen
til Dugdannelsen.

Denne Metode kan dog kun føre til en foreløbig Formodning.
Den henleder Opmærksomheden paa Forholdet mellem
$A$ og $B$, men udelukker ikke Muligheden af, at $B$ engang
kunde indtræde uden $A$. (Vi have kun Kombinationen $AB$ given
og vide Intet om de negative Tilfælde. Smlgn. 46.) Sker Fortætningen
ikke, naar pludselig Afkøling udelukkes? Dette
Spørgsmaal maa besvares, før Slutningen er sikker. Brugen
af et Lægemiddel kan have vist sig gavnlig i en Række Tilfælde;
men er det gavnligt i alle Tilfælde, og --- vilde Patienterne
ikke være blevne helbredede, selvom de ikke havde
brugt Lægemidlet? --- Vi føres herved over til den næste
Metode.

b. \textit{Forskelsmetoden}. Naar et Fænomen $B$ ikke indtræder,
uden Fænomenet $A$ gaar forud, medens alle andre Omstændigheder
ere ens (saa at $A$’s og $B$’s Indtræden udgør den eneste
Forskel), sluttes, at $A$ er Aarsag til $B$ (eller at $A$ og $B$ ere
Virkninger af samme Aarsag).
% scheme n2r-03-04 (notatex.py)
\[\begin{array}{l@{\qquad}l}
AVX & VX \\
BV_1X_1 & V_1X_1
\end{array}\]
Her behøves kun to Tilfælde, et, hvor $A$ og $B$ indtræde,
og et andet, hvor de under ellers lige Omstændigheder ikke
indtræde. --- Ved skydækket Himmel eller i Blæst sker ingen
Dugdannelse, kun i stille og klart Vejr. Ved at udspænde et
Lagen i nogle Fods Højde over Jorden hindrer man Varmeudstraalingen
fra denne, og da falder ogsaa Dugdannelsen bort
under Lagenet, medens den gaar for sig lige udenfor. --- Det
gælder om at finde en enkelt Omstændighed, som er afgørende
% --- p. 44 ---
% running head: Tillæg .
\opage{44}(idet den enten er hele Aarsagen, eller indeholder Aarsagen,
eller er en Del af Aarsagen, der altid bestaar i en hel Gruppe
af Betingelser, smlgn. Psykologi VD, 5 d; --- denne Bemærkning
gælder for alle fire Metoder, men er af særlig Vigtighed her).
Gray opdagede, at Forskellen mellem elektriske Ledere og
Isolatorer beroede paa Stoffernes Beskaffenhed; han havde
først antaget, at naar en Silkesnor var en bedre Isolator end
en Hampesnor, kom det af, at den var tyndere; men det viste
sig, at en Metaltraad, der var ligesaa tynd som Silkesnoren,
afledte Elektriciteten. --- Et Menneskes brødefri Vandel kan
skyldes hans gode Karakter --- men maaske ogsaa hans lykkelige
Forhold, der ingen Fristelser medføre. --- Efter de danske
Landboreformer i Slutningen af det 18. Aarhundrede indtraadte
en stor Prisstigning paa Landbrugsprodukter, saa at man ikke
klart kunde se, hvad der skyldtes dem; deres heldige Indflydelse
blev derimod tydelig, da Priserne i det 19. Aarhundrede
begyndte at falde igen. ---

Ved Forskelsmetoden have vi Kombinationerne $AB$ og $ab$
givne, hvoraf følger $A = B$ (der her er at fortolke saaledes:
hvis $A$ er, er $B$, og $B$ indtræder ikke uden $A$). (Smlgn. 17.)

De to sidste Metoder ere speciellere Former af Forskelsmetoden.


c. \textit{Restmetoden}. Naar $AVX$ gaar forud for $BV_1 X_1$, og naar
det vides, at $V$ er Aarsag til $V_1$ og $X$ til $X_1$, sluttes, at $A$ er
Aarsag til $B$.
% scheme n2r-03-05 (notatex.py)
\[\begin{array}{l@{\qquad}l@{\qquad}l}
AVX & V & X \\
\cdashline{1-3}
BV_1X_1 & V_1 & X_1
\end{array}\]
Af de Uregelmæssigheder i Planeten Uranus’ Bane, der
ikke kunde forklares ved de bekendte Planeters Tiltrækning,
sluttede man, at der udenfor Uranus maatte findes en hidtil
ukendt Planet, og dette ledte til Opdagelsen af Neptun. --- Menneskelige
Karakterer og Raceejendommeligheder kunne ikke blot
forklares ved Opdragelsens og andre ydre (fysiske og sociale)
Forholds Indflydelse; der bliver en uforklaret Rest tilbage, som
maa forklares ved Arvelighed. (Smlgn. Psykologi VII C, 2.)

d. \textit{De proportionale Forandringers Metode}. Naar et Fæ%
% --- p. 45 ---
% running head: 1. Mill's induktive Metoder .
\opage{45}nomen $B$ tiltager eller aftager, eller idethele forandrer sig, altefter
som et andet Fænomen $A$ undergaar tilsvarende Forandringer,
medens alle andre Omstændigheder forblive uforandrede,
sluttes, at $A$ er Aarsag til $B$ (eller at begge ere
Virkninger af samme Aarsag).
% scheme n2r-03-06 (notatex.py)
\[\begin{array}{l@{\qquad}l@{\qquad}l}
A_1VX & A_2VX & A_3VX ......... \\
\cdashline{1-3}
B_1V_1X & A_2V_1X_1 & B_3V_1X_1
\end{array}\]
% PRINTED AS IS: the middle cell of the bottom row is A₂V₁X₁ (B₂V₁X₁ expected?), and B₁V₁X has no subscript on X (z-06-04)
Denne Metode finder Anvendelse, hvor man ikke helt kan
udelukke et Fænomen. Man undersøger da de Tilfælde, hvor det
varierer, og søger at finde andre Fænomener, der under samme
Omstændigheder variere paa tilsvarende Maade. Metoden kræver,
at de varierende Rækker fortsættes saa langt som muligt.

Inertisætningen begrundedes af Galilei ved denne Metode:
jo glattere han gjorde sine Kugler og de Render, han lod dem
løbe i, des længere bleve Kuglerne ved at løbe. --- At Sjælelivet
særligt er knyttet til Nervesystemet (hvad enten det ene
er Aarsag til det andet, eller Aarsagen ligger længere tilbage),
sluttes af, at de variere paa tilsvarende Maade. (Psykologi
II, 6---7.) --- Statistiken viser, at Ægteskabernes Antal stiger, naar
Kornpriserne synke, og falder, naar disse stige. ---

Med Hensyn til disse Mill’ske Metoder maa mærkes, at de
forudsætte en tydelig Forskel mellem Aarsag og Virkning. En
saadan tydelig Forskel er der kun for den foreløbige Opfattelse;
naar Forskningen trænger længere ind i Forholdet,
finder den en kontinuerlig Overgang fra det, der kaldes „Aarsag“,
til det, der kaldes „Virkning“. (Smlgn. Psykologi VD, 3.)

De Mill’ske Metoder forudsætte Aarsagssætningens Gyldighed
og ligge forsaavidt udenfor den formelle Logik (smlgn. 22).
Mill indrømmede dette og kom derved til at bevæge sig i en
Kreds, da han begrundede selve Aarsagssætningen ved Induktion
(Overensstemmelsesmetoden). (Smlgn. „Oversigt over den
nyere Filosofis Historie“ VII b.)

Spørgsmaalet om Aarsagssætningens Gyldighed fører over
fra den formelle Logik til Erkendelsesteorien, der spørger, med
hvilken Ret vi tillægge Begrundelsens Princip (og de deri inde%
% --- p. 46 ---
% running head: Tillæg .
\opage{46}holdte logiske Principer, se 21 til 29) Gyldighed for den reale
Tilværelse, ikke blot for Ordningen af vore Tanker. --- Ogsaa
i Psykologien komme vi til at drøfte Aarsagsforholdet, nemlig
ved Spørgsmaalet om, hvorledes vi skelne mellem Virkelighed
og blot Forestilling (Illusion eller Drøm). (Smlgn. herom Psykologi
VD, 1---2.)

\sechead{2. Analogi og Analogislutning.}{\textbf{2. Analogi og Analogislutning.}}


Den inderste Nerve i al Logik er Identitet. Det er i Kraft
af den, at vi fæstne vore Forestillinger til Begreber, knytte
Begreber sammen i Domme og gaa fra Domme til Domme
gennem Slutning. Vor Tænknings Maal er derfor saavidt som
muligt at reducere Forskelligheder til Identiteter --- eller dog
at opfatte dem, som om de vare Identiteter, og se, hvor langt
vi kunne gennemføre denne Opfattelse. Men Tanken fører en
stadig Kamp med det Givne for at kunne naa dette Maal. Og
paa flere Punkter maa Analogi (Forholdslighed) træde i Stedet
for Identitet.

a. Vore Begreber bero for en stor Del, nemlig afset fra
„Mærkebegreberne“, paa Analogi mellem de Fænomener, for
hvilke de gælde (5). Kun tilnærmelsesvis kan der mellem forskellige,
i Anskuelsen givne Fænomeners Kendemærker paavises
en saa streng Identitet, som Identitetsprincipet kræver
(24). Der fremtræder her et irrationelt Forhold mellem Erkendelsens
Stof og Form, mellem det Givne og dets Udtryk i
Tanken. Have vi éngang dannet et Begreb, maa vi arbejde
med det som identisk med sig selv, indtil vi maaske af forøget
og klarere Erkendelsesstof kunne danne et fuldkomnere
Begreb. Men det maa ikke glemmes, at det ofte fra først af
beror paa Analogi, og at vore Definitioner ofte ere Hypoteser.

b. Naar vi nu forbinde Begreberne i Domme, og Dommene
i Slutninger, kunne de Forbindelser, der saaledes tilvejebringes
mellem vore Tanker, ikke uden videre antages at gengive de
Forbindelser, der finde Sted mellem de givne Fænomener.
Forholdet mellem Grund og Følge paa den ene Side og mellem
Aarsag og Virkning paa den anden Side have vi ikke Ret
til at identificere. Vi kunne opfatte Begivenhedernes Række%
% --- p. 47 ---
% running head: 2. Analogi og Analogislutning .
\opage{47}følge, som om den gengav et Forhold mellem Grund og Følge;
men kun en Analogi er det her berettiget at antage (22). Paa
lignende Maade forholder det sig ved Anvendelsen af matematiske
Begreber og Domme paa Fænomenerne. Talrækken er den
simpleste og mest overskuelige Ordning af Elementer, vi kende,
og vi søge derfor naturligt at ordne Fænomenerne i analoge
Rækker, saa at vi fra de Forhold, der gælde mellem Tallene,
kunne slutte til Forhold mellem Fænomenerne. Den saakaldte
mekaniske Naturopfattelse beror paa den store Analogi mellem
Bevægelser i Rummet og kvalitative Forandringers Rækker.
Saaledes betragter Fysiken Farverækken som om den blot var
en Række kvantitative Ændringer (i Æterbølgers Brydning
o. s. v.), og Stoffernes kvalitative Forandringer betragtes i Kemien,
som om de blot vare Omlejringer af uforandrede
Smaadele. Anvendelsen af de formelle Videnskaber paa de
reale Videnskaber (smlgn. 2 og 22) skyldes altsaa Analogi.

c. I de nys nævnte Exempler optræder Analogislutningen
i dens mest enkelte Form, som de proportionale Forandringers
Metode. Slutning fra den ene Række (f. Ex. den kvantitative)
kan her stedse finde sin Bekræftelse ved en Iagttagelse af nye
Led i den anden Række (den kvalitative).

Men det er ikke altid, at Analogislutning kan antage denne
Form. Ofte er den en Slutning i anden Figur (33): naar to
Subjekter have en fælles Egenskab, ligger det nær at spørge,
om de da ikke stemme overens i flere Egenskaber, især naar
den fælles Egenskab er af væsentlig, dybt indgribende Betydning.
At baade Fiske og Hvaler ere Vanddyr, kan foranledige
en Undersøgelse af, hvilke andre Egenskaber de have fælles.
Medens anden Figur ikke kan begrunde nogen logisk Slutning,
kan den give Stødet til en Opdagelse, der saa kan finde
sin Bekræftelse ad helt andre Veje. Og naar et Subjekt har
to Prædikater (tredie Figur, se 34), ligger det nær at spørge,
hvorledes disse to Prædikater (f. Ex. Vanddyr og Hvirveldyr
om Fiskene) indbyrdes bestemme hinanden; det kan lede til
at opsøge flere andre Tilfælde, i hvilke de to Prædikater forekomme
sammen. (Man gør da egentligt de to Prædikater til
Subjekter, saa at tredie Figur gaar over til anden Figur, der er
% --- p. 48 ---
% running head: Tillæg .
\opage{48}det egentlige Skema for Analogislutning.) Saaledes kan Analogien
virke fremaddrivende, og det paa en mere positiv Maade
end Selvmodsigelsen (se 46). Den er en betydningsfuld Opdagelsesmetode.
Men naar Beviset skal føres, maa det føres
ved Hjælp af Identiteter. Newtons Opdagelse af Tyngden som
Verdenslov skal være bleven gjort ved Hjælp af en dristig
Analogi (10) mellem Maanens Bevægelse og Faldbevægelsen
paa Jorden; men Beviset førtes efter Formlen $A = B = C$ (32),
hvor $B$ betyder de to Bevægelsers identiske Lov.

Træffende har Aristoteles, der kun bruger Ordet Analogi
i Betydning af matematisk Proportion, kaldt det, vi benævne
Analogislutning, for Slutning fra et Exempel (paradeigma). Vi
slutte fra Forholdene i et Tilfælde til Forholdene i et andet
Tilfælde. Den danner hos ham en Modsætning baade til Induktion
og til Deduktion. --- Den nærmer sig til Induktion, jo flere
forskellige Exempler den støtter sig til, saa at en Dom kan
begrundes og maaske senere bekræftes.

d. Det er ved al Analogiseren det store Spørgsmaal, om
vi have Ret til at overføre, hvad der gælder i én Henseende
eller for visse Fænomener, til andre Henseender eller andre
Fænomener, fordi et enkelt Lighedspunkt kan paavises. Tilsidst
stiller Spørgsmaalet sig, og det i sin mest problematiske Form,
naar vi staa ved Grænsen af vor Erkendelse og spørge, om
vi fra det, vor Erfaring og Tænkning lærer os om Tilværelsen,
kunne drage Slutninger til dennes hele Natur. Her har Menneskeheden
--- i sine Religioner og i sine Spekulationer
--- arbejdet med store og dristige Analogier, der ikke kunne bekræftes
ved Erfaring, og hvis Værd heller ikke behøver at
bero derpaa.

\end{document}
